%------------------------------------------------------------------------------%
% TOPOLOGICAL AND NORMED LINEAR SPACES                                         %
%------------------------------------------------------------------------------%
% File  : Topological_and_Normed_Linear_Spaces.tex                             %
% Author: Klaus Hoeflinger, Beethovenstrasse 11, D-73117 Wangen                %
% Email : khoeflinger@t-online.de                                              %
%------------------------------------------------------------------------------%

%------------------------------------------------------------------------------%
%                       DOCUMENT CLASS and INCLUDE FILES                       %
%------------------------------------------------------------------------------%
\documentclass[11pt,twoside]{article}
\usepackage{geometry}
\geometry{paperheight=241mm,
          paperwidth=152mm,
          top=22mm,
          bottom=17mm,
          left=20mm,
          right=20mm,
          textwidth=112mm,
          textheight=202mm,
          headheight=10mm,
          headsep=3mm,
          footskip=9mm,
          includehead,
          includefoot,
          marginparsep=0mm,
          marginparwidth=0mm}

\renewcommand{\baselinestretch}{0.945}\normalsize

\AtBeginDocument{\setlength\abovedisplayskip{2mm}}
\AtBeginDocument{\setlength\belowdisplayskip{2mm}}

%------------------------------------------------------------------------------%
% define title formats                                                         %
%------------------------------------------------------------------------------%
\usepackage{titlesec}

\renewcommand{\thesection}{\arabic{section}}
\renewcommand{\sfdefault}{FiraSans}
\renewcommand{\ttdefault}{pcr}
\renewcommand{\rmdefault}{antiqua}

\titleformat{\part}[display]
{\normalfont \large \rmfamily \bfseries \centering}
{\small{\thepart.}}{2pt}{}

\titleformat{\section}[runin]
{\normalfont \rmfamily \bfseries}
{\thesection}{1pt}{.\;}[.]

%\titleformat{\section}
%{\normalfont \bfseries \small \filcenter}
%{\thesection.\;}{0mm}{}

\titleformat{\subsection}[runin]
{\normalfont \itshape}
{}{0mm}{}

\titleformat{\subsubsection}[runin]
{\normalfont \itshape}
{}{}{}{}

\titleformat{\paragraph}[runin]
{\normalfont}
{}{}{}{}

\titleformat{\subparagraph}[runin]
{\normalfont}
{}{}{}{}

\titlespacing*{\part}          { 0pt}{ 0pt}{ 8pt}
\titlespacing*{\section}       { 0pt}{ 7pt}{ 4pt}
\titlespacing*{\subsection}    { 0pt}{14pt}{ 6pt}
\titlespacing*{\subsubsection} { 0pt}{ 6pt}{12pt}
\titlespacing*{\paragraph}     { 0pt}{ 6pt}{12pt}
\titlespacing*{\subparagraph}  { 0pt}{ 0pt}{12pt}

%------------------------------------------------------------------------------%
% define enumerate parameters                                                  %
%------------------------------------------------------------------------------%
\usepackage{enumitem}
\setlength{\parindent}{3pt}
\setlength{\parskip}  {0pt}

%\setlist{noitemsep}

\setlist[enumerate]{
         leftmargin=12pt,
         rightmargin=0pt,
         itemsep=1pt,
         itemindent=3pt,
         topsep=3pt,
         labelsep=4pt,
         partopsep=0pt,
         labelwidth=0pt,
         listparindent=0pt,
         parsep=0pt}

\setlist[itemize]{
         leftmargin=12pt,
         rightmargin=0pt,
         itemsep=0pt,
         itemindent=1pt,
         topsep=1pt,
         labelsep=4pt,
         partopsep=1pt,
         labelwidth=0pt,
         listparindent=0pt,
         parsep=0pt}

%------------------------------------------------------------------------------%
% define packages                                                              %
%------------------------------------------------------------------------------%
\usepackage{upref}
\usepackage{amssymb}
\usepackage{slantsc}
\usepackage{amsmath}
\usepackage{amsthm}
\usepackage{thmtools}
\usepackage{amscd}
\usepackage{verbatim}
\usepackage{latexsym}
\usepackage{multicol}
\usepackage{graphicx}
\usepackage{setspace}
\usepackage[ngerman,english]{babel}
\usepackage[protrusion=true,expansion=true]{microtype}
\usepackage{tikz}
\usetikzlibrary{arrows,arrows.meta,decorations.markings}
\usepackage{mathrsfs}
\usepackage{amsfonts}
\usepackage{datetime2}
\usepackage{textgreek}
\usepackage{hyperref}
\usepackage{pifont}
\usepackage{relsize}
%\usepackage{biblatex}
\relsize{-0.05}
\usepackage[skip=0mm]{parskip}
\usetikzlibrary{decorations.pathmorphing}

\usepackage{mlmodern}
\usepackage[T5,T1]{fontenc}
\usepackage{ragged2e}

%\usepackage{mathabx}
%\usepackage{times}

%\usepackage{fontspec}
%\setmainfont{Nimbus Roman No9 L}
%\usepackage[T1]{fontenc}

%------------------------------------------------------------------------------%
% define page layout                                                           %
%------------------------------------------------------------------------------%
\usepackage{fancyhdr}
\pagestyle{fancy}
\setlength{\headheight}{30.0pt}
\renewcommand{\headrulewidth}{0pt}

\AtBeginDocument{
  \addtolength\abovedisplayskip{-0.0\baselineskip}
  \addtolength\belowdisplayskip{-0.0\baselineskip}
}

\fancyhead{}
\fancyfoot{}
\addtolength{\headwidth}{0.02\marginparwidth}

\makeatletter
\let\ps@plain\ps@fancy
\makeatother

\renewcommand\sectionmark[1]
{
 \def\sectionname{#1}
 \markright{\thesection #1}
}

\makeatletter
\@addtoreset{section}{part}
\makeatother

%------------------------------------------------------------------------------%
% define footnote without number                                               %
%------------------------------------------------------------------------------%
\newcommand{\myfootnote}[1]{%
\renewcommand{\thefootnote}{}%
\footnotetext{#1}%
\renewcommand{\thefootnote}{\arabic{footnote}}%
}

%------------------------------------------------------------------------------%
% define numbering in equations                                                %
%------------------------------------------------------------------------------%
\numberwithin{equation}{section}

%------------------------------------------------------------------------------%
% define newtheorem                                                            %
%------------------------------------------------------------------------------%
\declaretheoremstyle[
headfont=\scshape, 
headindent = 0mm, %\parindent,
%postheadhook = {\hspace{0mm}\newline},
%postheadspace = \newline, 
spaceabove = 3mm, 
spacebelow = 3mm,
numberwithin=section,
name=Theorem]{khMATH}
\declaretheorem[style=khMATH]{theorem}

\declaretheoremstyle[
headfont=\bfseries, 
headindent = 0mm, %\parindent,
%postheadhook = {\hspace{0mm}\newline},
%postheadspace = \newline, 
spaceabove = 3mm, 
spacebelow = 3mm,
numberwithin=part,
name=Theorem]{khMATH1}
\declaretheorem[style=khMATH1]{theorem1}

\declaretheoremstyle[
headfont=\scshape, 
headindent = 0mm, %\parindent,
%postheadhook = {\hspace{0mm}\newline},
%postheadspace = \newline, 
spaceabove = 3mm, 
spacebelow = 3mm,
numberwithin=section,
name=Corollary]{khMATH1}
\declaretheorem[style=khMATH1]{corollary}

%            name         counter     label              numeration-mode
\theoremstyle{plain}
\newtheorem*{introduction}            {\sc Introduction}
\newtheorem {standard}                {}                 [section]
\newtheorem*{definition}              {\sc Definition}
\newtheorem{satz1}                    {\sc Satz}
\newtheorem{lemma}                    {\sc Lemma}
\newtheorem*{proposition}             {\sc Proposition}
%\newtheorem{corollary}               {\sc Corollary}
\newtheorem*{corollar}                {\sc Korollar}
\newtheorem*{remark}                  {\sc Remark}
\newtheorem*{remarks}                 {\sc Remarks}
\newtheorem*{example}                 {Example}
\newtheorem*{examples}                {Examples}
\newtheorem*{theo}                    {\sc Theorem}

%------------------------------------------------------------------------------%
% define tabular                                                               %
%------------------------------------------------------------------------------%
\usepackage{tabularx}
\newcolumntype{L}[1]{>{\raggedright\arraybackslash}p{#1}}
\newcolumntype{C}[1]{>{\centering\arraybackslash}  p{#1}} 
\newcolumntype{R}[1]{>{\raggedleft\arraybackslash} p{#1}} 

%------------------------------------------------------------------------------%
% define \sh, \zB                                                              %
%------------------------------------------------------------------------------%
\def\dsh{{d.\,h.}}
\def\ie{{i.\,e.}}
\def\zB{z.\,B.}
\renewcommand{\abstractname}{ }

%------------------------------------------------------------------------------%
% change default spacing for cases                                             %
%------------------------------------------------------------------------------%
\makeatletter
\renewcommand*\env@cases[1][1.6]{%
  \let\@ifnextchar\new@ifnextchar
  \left\lbrace
  \def\arraystretch{#1}%
  \array{@{}l@{\quad}l@{}}%
}
\makeatother

\newenvironment{rcases}
  {\left.\begin{aligned}}
  {\end{aligned}\right\rbrace}

%------------------------------------------------------------------------------%
% change default for \predisplaypenalty                                        %
%------------------------------------------------------------------------------%
\predisplaypenalty=990
\allowdisplaybreaks \numberwithin{equation}{section}

%------------------------------------------------------------------------------%
% general notations                                                            %
%------------------------------------------------------------------------------%
\def\*{\star}
\def\={\,=\,}
\def\+{\,+\,}
\def\xsubset{{\,\subset\,}}
\def\xtimes{{\,\times\,}}
\def\xeof{{\,\in\,}}
\def\eq{{\,=\,}}
\def\xto{{\,\to\,}}
\def\eqdef{\,:=\,}
\def\:{{\,:\,}}
\def\ele{{\,\in\,}}
\def\exists{\mbox{ exists }}
\def\and{\mbox{ and }}
\def\for{\mbox{ for }}
\def\fa{\mbox{ for all }}
\def\fe{\mbox{ for each }}
\def\qfa{\quad \mbox{ for all }}
\def\df{\,\dotfill}
\def\iff{if and only if }
\def\wrt{with respect to}

\makeatletter
\newcommand \Df {\leavevmode \cleaders \hb@xt@ .70em{\hss .\hss }\hfill \kern \z@}
\makeatother

\def\sql{\mathfrak{a}}               % sesquilinear form a
\def\DF{B}                           % Dirichlet form a
%------------------------------------------------------------------------------%
% number sets and special elements                                             %
%------------------------------------------------------------------------------%
\def\P{\mathbb{H}}                   % half plane of $\C$
\def\J{I}                            % interval I
\def\N{{\mathbb{N}}}                 % unsigned integers
\def\Q{{\mathbb{Q}}}                 % rational integers
\def\Z{{\mathbb{Z}}}                 % integers
\def\R{{\mathbb{R}}}                 % real numbers
\def\Rn{{\mathbb{R}}^n}              % real numbers in Rn
\def\Rd{{\mathbb{R}}^d}              % real numbers in Rn
\def\C{{\mathbb{C}}}                 % complex numbers
\def\F{{\mathbb{K}}}                 % arbitrary field
\def\Torus{{\mathbb{T}}}             % complex circle line
\def\T{\tau}                         % upper bound of interval (0,t)
\def\sign#1{\operatorname{sign}(#1)} % sign of a number

%------------------------------------------------------------------------------%
% set theory                                                                   %
%------------------------------------------------------------------------------%
\def\G{G}                            % domain $G$
\def\clG{\overline{\G}}              % closure of $G$
\def\bndG{\partial \G}               % boundary of $G$
\def\setA{\mathcal{A}}               % set A
\def\setB{\mathcal{B}}               % set B
\def\setC{\mathcal{C}}               % set C
\def\setM{M}                         % set M
\def\setN{N}                         % set N
\def\setS{\mathcal{S}}               % set S
\def\famF{\mathfrak{F}}              % family of sets
\def\indI{\mathcal{I}}               % index set I
\def\pot#1{\mathfrak{P}(#1)}         % power set

\def\relR{\mathbf{R}}                % relation R
\def\mapf{f}                         % mapping f
\def\mapg{g}                         % mapping g

\def\set#1#2{\{ \, #1 \,: \, #2 \, \}}
\def\tensor#1#2{#1 \otimes #2}
\def\Tens{\oplus}

\def\cal#1{{\mathcal{#1}}}
\def\aut#1{\operatorname{Aut}(#1)}

\def\cross#1#2{#1 {\,\times\,} #2}
\def\multicross#1#2{#1 \times  \ldots \times  #2}
\def\multitens#1#2 {#1 \otimes \ldots \otimes #2}

%------------------------------------------------------------------------------%
% group theory                                                                 %
%------------------------------------------------------------------------------%
\def\1{{\mathsf{1}}}                % unit element of a monoid
\def\0{{\mathsf{0}}}                % unit element of an Abelian monoid
\def\kernel#1{\mathfrak{N}(#1)}     % kernel of a homomorphism
\def\cokernel#1{\mathfrak{C}{(#1)}} % cokernel of a homomorphism
\def\Hom#1{\operatorname{Hom}(#1)}

%------------------------------------------------------------------------------%
% linear algebra                                                               %
%------------------------------------------------------------------------------%
\def\vecV{\mathit{V}}               % vector space V
\def\vecH{\mathit{H}}               % vector space H
\def\vecW{\mathit{W}}               % vector space W
\def\vecU{\mathit{U}}               % subspace U
\def\convC{\mathcal{C}}             % convex set C
\def\basH{\mathfrak{B}}             % Hamel base B

\def\LO#1{\mathscr{L}(#1) }         % algebra of linear transformations

\def\scalar#1#2{ \langle #1,#2 \rangle }
\def\trace#1{\operatorname{tr}(#1)}
\def\dim#1{{\operatorname{dim} \, #1 }}
\def\codim#1{{\operatorname{codim} \, #1 }}
\def\dim{\operatorname{dim}}
\def\codim{\operatorname{codim}}
\def\ind{\operatorname{ind}}

\def\genG{\cal{G}}

\def\algA{\cal{A}}
\def\algB{\cal{B}}
\def\algU{\cal{U}}

\def\idI{\cal{I}}

%------------------------------------------------------------------------------%
% topology                                                                     %
%------------------------------------------------------------------------------%
\def\compK{{K}}          % compact set C
\def\openO{\mathcal{O}}  % open set O
\def\openU{\mathcal{U}}  % open set U
\def\U{\mathcal{U}}      % neighbourhood U
\def\V{\mathcal{V}}      % neighbourhood V
\def\d{\mathit{d}}       % metric function d

\def\interior#1{\mathring{#1}}
\def\closure#1{{\overline{#1}}}

\def\topT{\mathcal{T}}   % topology T
\def\topX{X}             % topological space X
\def\topY{Y}             % topological space Y
\def\topZ{Z}             % topological space Z
\def\metrS{M}            % metric space M

\def\grpG{G}             % topological group G
\def\grpA{A}             % topological group A
\def\grpB{B}             % topological group B
\def\grpC{C}             % topological group C
\def\grpH{H}             % topological group H
\def\grpU{U}             % topological group U
\def\grpV{V}             % topological group V
\def\topG{G}             % topological group G
\def\topH{H}             % topological group H
\def\topU{U}             % topological subgroup U

\def\subS{\cal{S}}
\def\riR{R}              % ring R

%------------------------------------------------------------------------------%
% calculus                                                                     %
%------------------------------------------------------------------------------%
\def\supp#1{\operatorname{supp}(#1)}
\def\singsupp#1{\operatorname{sing \, supp}(#1)}

\def\re{\operatorname{Re}}
\def\im{\operatorname{Im}}
\def\grad{\operatorname{grad}}

%\def\abs#1 {\left|#1\right|}
\def\abs#1 {\lvert #1 \rvert}
%\def\abs#1 {\left|\,#1\,\right|}

%------------------------------------------------------------------------------%
% measure theory                                                               %
%------------------------------------------------------------------------------%
\def\salgA{\Sigma}               % sigma-algebra S
\def\salgB{\cal{B}}              % sigma-algebra B
\def\Borel#1{\cal{B}_{#1}}       % Borel-algebra 
\def\LB{\lambda}                 % Lebesgue-Borel measure

%------------------------------------------------------------------------------%
% function spaces                                                              %
%------------------------------------------------------------------------------%
\def\Lp{L^p(\G)}                 % spaces L_p
\def\Lq{L^q(\G)}                 % spaces L_q
\def\Lh{L^2(\G)}                 % spaces L_2
\def\Li{L^{\infty}(\G)}          % spaces L_infty
\def\Lc{L_{loc}^1(\G)}           % spaces L_1^{loc}
\def\Ck{C^k(\G)}                 % spaces C_k

\def\BV#1{BV(#1) }               % set of functions of bounded variation
\def\AC#1{AC(#1) }               % set of absolutely continuous functions

\def\MP#1{\mathscr{M}^+(#1) }    % set of positive measures
\def\MS#1{\mathscr{M}^-(#1) }    % vector space of finite signed measures
\def\MC#1{\mathscr{M}(#1) }      % vector space of complex measures
\def\MCR#1{\mathscr{M}_{r}(#1) } % vector space of regular complex measures

\def\SobW{W^{k,p}(\G)}                % Sobolev space W
\def\SobO{\mathring{W}^{k,p}(\G)}     % Sobolev space W
%\def\WN#1#2{\mathring{W}^{#1}(#2) }   % Sobolev space W
\def\WN#1#2{W_0^{#1}(#2) }   % Sobolev space W
%\def\HN#1#2{\mathring{H}^{#1}(#2) }   % Sobolev space H
\def\HN#1#2{H_0^{#1}(#2) }   % Sobolev space H
%------------------------------------------------------------------------------%
% functional analysis                                                          %
%------------------------------------------------------------------------------%
\def\snorm#1{\lVert #1 \rVert}
\newcommand{\norm}[1]{\left\lVert #1 \right\rVert}

\def\topV{V}                       % topological vector space V
\def\topW{W}                       % topological vector space W
\def\normS{X}                      % normed space

\def\banX{\mathit{X}}              % Banach space X
\def\banY{\mathit{Y}}              % Banach space Y
\def\banZ{\mathit{Z}}              % Banach space Z
\def\dbanX{\mathit{X^*}}           % dual space of Banach space X
\def\dbanY{\mathit{Y^*}}           % dual space of Banach space Y
\def\dbanZ{\mathit{Z^*}}           % dual space of Banach space Z

\def\banA{\mathit{A}}              % Banach algebra A
\def\banB{\mathit{B}}              % Banach algebra B
\def\banC{\mathit{C}}              % C*-algebra C


\def\domain#1{\mathfrak{D}(#1)}    % domain
\def\range#1{\mathfrak{R}(#1)}     % range

\def\isoJ{\mathfrak{J}}            % isometry J
\def\repA{{\cal{A}_{\sql}}}        % representation operator A
\def\linS{{S}}                     % linear operator S
\def\linG{{G}}                     % linear operator G
\def\linL{{L}}                     % linear operator L
\def\linT{{T}}                     % linear operator A
\def\linA{\mathit{A}}              % linear operator A
\def\linB{{B}}                     % linear operator B
\def\linM{{M}}                     % linear operator M
\def\linU{{U}}                     % linear operator U
\def\linI{{I}}                     % linear operator I
\def\A{\cal{A}}                    % linear differential operator A
\def\BDO{\cal{B}}                  % linear boundary differential operator B
\def\linU{{U}}                     % unitary operator U
\def\linP{{P}}                     % projection P
\def\compA{k}                      % compact operator k
\def\linFR{f}                      % finite rank operator f
\def\fredA{A}                      % Fredholm operator F

\def\HAM{\cal{H}}                  % Hamilton operator
\def\PA{\partial^{\alpha}}         % elementary differential operator
\def\DO{\cal{A}}                   % linear differential operator
\def\BC{\cal{B}}                   % boundary operator
\def\SLO{\cal{L}_{p,q}}            % Sturm-Liouville operator
\def\MO{\cal{L}_{M}}               % momentum operator

\def\HS#1{S(#1)           }        % algebra of Hilbert-Schmidt operators
\def\FR#1{F(#1)           }        % algebra of finite rank operators
\def\TC#1{N(#1)           }        % algebra of trace class operators
\def\BU#1{\mathscr{U}(#1) }        % algebra of unitary linear operators
\def\BO#1{\mathscr{L}(#1) }        % algebra of bounded linear operators
\def\CO#1{\mathscr{C}(#1) }        % algebra of compact linear operators
\def\CL#1{\mathscr{C}(#1) }        % algebra of compact linear operators
\def\FO#1{\Phi(#1) }               % set of Fredholm operators
\def\FK#1#2{\Phi_{#2}(#1) }        % set of Fredholm operators with index k

\def\hilH{\mathit{H}}              % Hilbert space H
\def\clU{\mathit{U}}               % closed subspace U

\def\orthO{\mathfrak{O}}           % orthonormal system O
\def\orthS{\mathfrak{S}}           % orthonormal system S
\def\basB{\mathfrak{B}}            % basis B of a Hilbert space

\def\GL#1 {\mathbf{GL}(#1)}        % linear group
\def\OG#1 {\mathbf{O}(#1)}         % linear group
\def\SOG#1 {\mathbf{SO}(#1)}       % linear group

%------------------------------------------------------------------------------%
% distribution theory                                                          %
%------------------------------------------------------------------------------%
\def\disF{F}                  % distribution F
\def\disG{G}                  % distribution G
\def\disH{H}                  % distribution H
\def\disU{U}                  % distribution U
\def\disT{T}                  % distribution T
\def\SF{C^{\infty}(\G)}       % space of smooth functions
\def\TF{C_c^{\infty}(\G)}     % space of compactly supported smooth functions
\def\TFRd{C_c^{\infty}(\R^d)} % space of test functions
\def\SRd{\mathscr{S}(\R^d)}   % Schwartz space
\def\SR1{\mathscr{S}(\R)}     % Schwartz space
\def\B1#1{C^1(#1) }           % space of continuously differentiable functions
\def\Bm#1{C^m(#1) }           % space of continuously differentiable functions
\def\Ck{C^k(\G)}
\def\TD{\mathscr{S}'(\R^d)}   % space of temperated distributions
\def\E{\cal{E}}               % space of distributions with compact support
\def\D{\mathscr{D}}           % D
\def\FT{\mathcal{F}}          % Fourier transformation F
\def\FT{\mathscr{F}}          % Fourier transformation F

%------------------------------------------------------------------------------%
% other definitions                                                            %
%------------------------------------------------------------------------------%
\def\Proof{\textsc{Proof}}
\def\FRECHET{Fr\'{e}chet}
\def\ARZELA{Arzel\'{a}}
\def\GARDING{G\r{a}rding}
\def\HOELDER{H\"older}
\def\SCHROEDINGER{Schr\"odinger}
\def\JOERGENS{J\"orgens}
\def\WUEST{W\"uest}
\def\KAEHLER{K\"aehler}
\def\HOEFLINGER{H\"oflinger}
\def\POINCARE{Poincar\'{e}}
\def\FA{Functional Analysis}

\def\Author   {K.\,{\HOEFLINGER}}
\def\AuthorUC {K.\,HOEFLINGER}
\def\Version  {1.0}
\def\Email    {khoeflinger@t-online.de}


%\renewcommand\thesection{\Roman{section}}

\def\Title{Topological and Normed Linear Spaces}

\titleformat{\section}
{\normalfont \bfseries \filcenter}
{\thesection.\;}{0mm}{}

\titleformat{\subsection}[runin]
{\normalfont \filcenter}
{}{0mm}{}

\titleformat{\subsubsection}[runin]
{\itshape \filcenter}
{}{0mm}{}

\titlespacing*{\section}      {5pt}{12pt}{ 4pt}
\titlespacing*{\subsection}   {0pt}{ 6pt}{12pt}
\titlespacing*{\subsubsection}{0pt}{ 6pt}{ 3pt}

%------------------------------------------------------------------------------%
%                                BEGIN OF DOCUMENT                             %
%------------------------------------------------------------------------------%
\begin{document}
\thispagestyle{empty}

\centerline{\large{\textbf{\Title}}}
\vspace{4pt}
\centerline{\small{{\textsc{\Author}}}}
\vspace{10pt}

\fancyhead[C] {\small{\textsc{\Title}}}
\fancyhead[OL]{\small{\thepage}}
\fancyhead[ER]{\small{\thepage}}

\myfootnote
{
  Version {\Version} from \today;
  Email:  \textit{khoeflinger@\,t-online.de}
}

%------------------------------------------------------------------------------%
%                                 INTRODUCTION                                 %
%------------------------------------------------------------------------------%
It is our aim in this article
to review those fundamental classes of topological linear spaces,
in which the analysis of so-called "linear operators" is usually done.
These spaces include,
besides some generic classes like
{locally convex spaces} and {completely metrizable spaces},
especially the class of {normed linear spaces} and related subclasses,
where we want to emphasize
{Banach} and {Hilbert spaces}.

%------------------------------------------------------------------------------%
\section{Topological Linear Spaces}                                            %
%------------------------------------------------------------------------------%
Basically,
the scalar field $\F$ of the linear spaces under consideration
will either be the field $\R$ of real numbers
or the field $\C$ of complex numbers.
Notice that these fields are furnished
with the so-called {standard topology},
whose subbase is given by the totality of open intervals in the real line
and by the totality of open disks in the complex plane,
respectively.
In general,
if results on topological or normed linear spaces
prove to be independent of the choice of the scalar field,
we shall use the symbol $\F$ instead of $\R$ or $\C$. 

\subparagraph*{}
By a \textit{topological linear space} is meant an ordinary vector space $\vecV$
over a topological field $\F$ (usually that of real or complex numbers),
together with a topology $\topT_{\,\vecV}$ 
such that addition $(x,y) \mapsto x\,{+}\,y$
and outer multiplication $(\lambda,x) \mapsto \lambda x$ are continuous.
Hence a multiplicity of topological concepts,
like
the \textit{closure} of a set,
the \textit{neighbourhood} of a point,
the \textit{convergence} of a sequence,
the \textit{continuity} of a mapping and others
are meaningful in the context of topological linear spaces.

\paragraph*{} %- Hausdorff's separation axiom
Let us continue with the introduction of a basic property
that our topological linear spaces under consideration have to own:
all singletons $\{x\}$ with $x \in \vecV$ are closed sets,
which is equivalent to say
that the topology $\topT_{\,\vecV}$
fulfills \textsc{Hausdorff}'s separation axiom.
In other words,
each pair of points $x,y \in \banX$ can be separated
by open sets $\openO_x, \openO_y$,
that is,
there are $\openO_x \subset \vecV$ and $\openO_y \subset \vecV$ with
$x \in \openO_x$,
$y \in \openO_y$ and
$\openO_x \cap \openO_y \eq \emptyset$.
The validness of this separation axiom,
in general topology also called the \textit{T2-separation axiom},
implies not only
that every singleton $\{x\}$ is closed,
but also
that the limit of a convergent sequence of elements in $\vecV$ is unique.

\paragraph*{}
Among a lot of properties,
a topological linear space may have,
there are two special ones
that are of importance for our purposes:

\begin{itemize}[leftmargin=15pt,itemsep=3pt]
\item[1.]
A topological linear space $\vecV$ is said to be \textit{metrizable},
if there is a function
$\d\:\vecV \times \vecV \xto [0,\infty)$
fulfilling the properties of a \textit{metric}\footnote{
\,Given a non-empty set $\metrS$,
a mapping $\d\:\cross{\metrS}{\metrS} \xto \R$ is called a \textit{metric}
on $\metrS$
iff the following four conditions hold for all $x,y,z \in \metrS$:
\begin{itemize}[leftmargin=40pt,itemsep=0pt]
\item[(M-1)]
$d(x,y) \geq 0$,
\item[(M-2)]
$d(x,y) \eq 0$ if and only if $x=y$,
\item[(M-3)]
$d(x,y) \eq d(y,x)$;
\item[(M-4)]
$d(x,z) \leq d(x,y) + d(y,z)$ \hspace{3mm}(\textit{triangle inequality}).
\end{itemize}
Such a function $\d$ gives rise to a topology in $\metrS$,
called the \textit{metric topology} associated with $\d$,
and a subbase is given by all open balls
$B_r(x) \eqdef \{y \in \metrS: \d(x,y) < r\}$
with center $x \in \metrS$ and radius $r > 0$.
}
such that the metric topology associated with $\d$
resembles the given topology $\topT_{\,\vecV}$.

\item[2.]
A sequence $\{x_n\}_{n \in \N}$ with values
in a topological linear space $\vecV$ is said to be Cauchy
if for each neighbourhood $\U$ of $\0 \in \vecV$
there is an integer $N_{\U}$
such that $x_n - x_m \in \U$ for all $m,n > N_{\U}$.
Then $\vecV$ is called \textit{sequentially complete},
if every Cauchy sequence converges to a unique element $x \in \vecV$.
In the setting of metric spaces,
however,
the definition of Cauchy sequence may be simplified by means of the metric:
$\{x_n\}_{n \in \N}$ is Cauchy,
if for each $\epsilon > 0$ there is an integer $N_{\epsilon}$ 
such that
$x_m - x_n \in B_{\epsilon}(\0)$ for all $m,n > N_{\epsilon}$,
where $B_{\epsilon}(\0) \eqdef \{x \in \vecV: \d(x,0) < \epsilon\}$
is the open ball with radius $\epsilon$ and center $\0 \in \vecV$.
\end{itemize}

\subparagraph*{}
For a topological linear space being metrizable and sequentially complete
at the same time,
both concepts of completeness are compatible to each other,
that is,
such a space is complete as a metric space
\iff
it is sequentially complete as a topological linear space.
This equivalence is founded by the fact
that a metric space fulfills the \textit{first countability axiom}.

\subparagraph*{} %- continuous embeddings
In the sequel we also use the notion of \textit{continuous embedding}
of a topological linear space $\banX$ into a topological linear space $\banY$,
which is nothing but a continuous and \textit{injective} linear mapping
$\iota\:\banX \xto \banY$.
The linearity of $\iota$ means that 
$\iota(\alpha x + \beta y) \eq \alpha \iota(x) + \beta \iota(y)$ holds
for all $\alpha, \beta \in \F$ and $x,y \in \banX$.

\subparagraph*{} %- linear homeomorphisms
In addition,
if $\iota$ is surjective and $\iota^{-1}$ is continuous as well,
then both mappings are called \textit{linear homeomorphisms}
between $\banX$ and $\banY$.
By $\banX \cong \banY$ we mean
that $\banX$ and $\banY$ are linear-homeomorphic to each another.

%------------------------------------------------------------------------------%
\section{Families of Seminorms}                                                %
%------------------------------------------------------------------------------%
The class of \textit{locally convex spaces} can be seen as the most important
one within all kinds of topological linear spaces.
The topology $\topT_{\,\vecV}$ of a locally convex topological linear space
distinguishes by the fact
that the origin $\0$ in $\vecV$
(and thus any other element) owns a neighbourhood base
consisting of \textit{convex}, \textit{balanced} and \textit{absorbing} sets.

\subparagraph*{} %- seminorms
It turns out
that each such pair $(\vecV,\topT_{\,\vecV})$ gives rise to
a family of \textit{seminorms} on the space $\vecV$.
By the concept of seminorm
we mean a mapping $\abs{\,\cdot\,} \:\vecV \xto \R$ having the properties
\begin{itemize}[leftmargin=55pt,itemsep=2pt]
\item[(SN-1)] \quad
$\abs{\lambda x} \= \abs{\lambda} \cdot \abs{x} \for \lambda \in \F$
\textit{(homogenity)},
\item[(SN-2)] \quad
$\abs{x+y} \, \leq \, \abs{x} + \abs{x} \for x,y \in \vecV$
\textit{(subadditivity)},
\end{itemize}
which especially implies $\abs{x} \leq 0$ for $x \in \vecV$
and $\abs{\0} \eq 0$.

\subparagraph*{} %- Minkowski functional
Moreover,
if $\U \subset \vecV$ forms a convex, balanced and absorbing neighbourhood
of the origin $\0$,
then the \textit{Minkowski functional}
\[
\begin{cases}[1.3]
\quad \abs{\,\cdot\,} _{\,\U}\:\vecV \xto \R_0^+,\\
\quad x \mapsto \abs{x} _{\,\U} \, \eqdef \,
\mbox{inf} \, \set{\lambda \geq 0}{ x \in \lambda \, \U} \\
\end{cases}
\]
fulfills all the properties of a seminorm.
On the other hand,
a family $\famF$ of seminorms defined on a linear space $\vecV$
leads in a natural way to a locally convex topology $\topT_{\,\vecV}$:
it is the coarsest topology on $\vecV$
such that all seminorms in $\famF$ remain continuous.

\subparagraph*{}
In summary,
there is a one-to-one correspondence
between linear spaces endowed with locally convex topologies
and those endowed with families of seminorms.
Notice that a locally convex topology is Hausdorff
iff
the family of seminorms is \textit{separating},
{\ie} for each element $x \neq 0$
there is at least one seminorm $\abs{\,\cdot\,} _{\alpha}$
such that $\abs{x} _{\alpha} \neq 0$.

\paragraph*{} %- F-spaces and Frechet spaces
There are a two subclasses of locally convex spaces,
which in the following we want to review:

\begin{itemize}[leftmargin=15pt]
\item[1.]
A topological linear space $\vecV$ is called \textit{F-space},
if it is also a \textit{complete} metric space with translation-invariant
metric $\d$,
{\ie}
\[
\d(x,y) \= \d(x+z,y+z) \; \fa x,y,z \in \vecV,
\]
and the metric topology coincides with the given topology on $\vecV$.

%\subparagraph*{} %- Frechet space
\item[2.]
If the topology of the F-space $\vecV$ is locally convex,
then $\vecV$ is said to be a \textit{{\FRECHET} space}\footnote{
A metrizable locally convex space that is not necessarily complete,
is called a \textit{pre-{\FRECHET} space}.}.
Alternatively,
a locally convex space is {\FRECHET}
\iff
it is complete and its family of seminorms is mostly countable.
In this case
we may define a translation-invariant metric on a {\FRECHET} space $\vecV$
according to
\[
\d(x,y) \, \eqdef \,
\sum_{n=1}^{\infty} p_n \cdot \frac{\abs{x-y} _n}{1+ \abs{x-y} _n}
\quad (x,y \in \topV)
\]
where $\{p_n\}_{n=1}^{\infty}$ is any sequence of positive real numbers,
whose infinite sum converges.
\end{itemize}

\paragraph*{} %- continuous linear functionals and adjoint spaces
Given a topological $\F$-linear space $\banX$,
let $\banX^*$ be the set of all continuous linear mappings
$\phi\:\banX \xto \F$,
called the \textit{adjoint space} $\banX^*$ of $\banX$,
which forms a linear space again.
If the topology of $\banX$ is at least locally convex,
then a specific version of the celebrated \textsc{Hahn-Banach} \textit{theorem}
ensures that
there are non-trivial continuous linear mappings on $\banX$,
thus we may assume $\banX^* \neq \{\0\}$.
We denote the elements of the adjoint space $\banX^*$ as
\textit{continuous linear functionals}.

\subparagraph*{} %- bi-adjoint spaces and weak*-topologies
It is known from linear algebra
that an ordinary $\F$-linear space $\vecV$ may be regarded
as a linear subspace of its (algebraic) bi-adjoint space
$\vecV^{\#\#} \eqdef (\vecV^{\#})^{\#}$,
where $\vecV^{\#}$ is the space of linear functionals $\phi\:\vecV \xto \F$.
More precisely,
every element $x \in \vecV$ gives rise to a linear functional
$F_x\:\vecV^{\#} \xto \F$ according to
$F_x(\phi) \eqdef \phi(x)$ for all $\phi \in \vecV^{\#}$.
The natural mapping $\iota\:\vecV \xto \vecV^{\#\#}$,
defined by $\iota(x) \eqdef F_x$ for $x \in \vecV$,
is linear and injective but in general not surjective.

\subparagraph*{} %- weak*-topologies
Having these facts in mind,
we introduce the \textit{weak*}-topology on the adjoint $\banX^*$
of a locally convex topological linear space $\banX$:
it is the smallest topology in $\banX^*$
such that
for each $x \in \banX$ the linear functional $F_x$ remains continuous.

%------------------------------------------------------------------------------%
\section{Normed Linear Spaces}                                                 %
%------------------------------------------------------------------------------%
Based on the results obtained in the last section,
a normed linear space is nothing but a pre-{\FRECHET} space,
whose family of separating seminorms consits of only one element
$\norm{\,\cdot\,} \eqdef \abs{\,\cdot\,} _0$,
which of course must be a separating one,
that is,
$\norm{x} \eq 0$ implies $x \eq \0$.

\subparagraph*{}
It is possible to introduce the concepts of norm and normed linear space
in a direct way, 
{\ie} without using the theory on families of seminorms
and locally convex spaces.
Following this approach,
a \textit{normed space} $\banX$ is a $\F$-linear space
that differs from other abstract linear spaces
by the presence of a particular mapping $\norm{\,\cdot\,}: \banX \xto \R$,
denoted the \textit{norm} on $\banX$,
which fulfills for all $x,y \in \banX$ and $\lambda \in \F$

\begin{itemize}[leftmargin=55pt,itemsep=2pt]
\item[(N-1)] \quad
$\norm{\lambda \, x} \= \abs{\lambda} \cdot \norm{x}$,

\item[(N-2)] \quad
$\norm{x+y} \, \leq \, \norm{x} + \norm{y}$
\; \textit{(triangle inequality)},

\item[(N-3)] \quad
$\norm{x} \geq 0$ and $\norm{x} \eq 0 \Leftrightarrow x \eq \0$.
\end{itemize}

\paragraph*{}
Linear spaces can be furnished with different norms.
Then two arbitrary norms $\norm{\,\cdot\,}_{\alpha}$
and $\norm{\,\cdot\,}_{\beta}$ defined on a linear space $\banX$
are called \textit{equivalent} to each another,
if there are constants $c_1,c_2 > 0$ such that
$
c_1 \norm{x}_{\alpha} \leq \norm{x}_{\beta}  \leq c_2 \norm{x}_{\alpha} 
$
for all $x \in \banX$.

\subparagraph*{} %- norm topologies
Each normed linear space $\banX$ can be equipped
with a very natural topology $\topT$,
called the \textit{norm topology} on $\banX$.
It can be defined either by the assertion
that $\topT \eqdef \topT_{\norm{\cdot}}$ is the smallest topology on $\banX$
such that the norm mapping remains continuous,
or it is defined to be the metric topology $\topT_d$,
when $\banX$ is considered as a metric space with metric function
$\d(x,y) \eqdef \norm{x-y}$ for all $x,y \in \banX$.
In fact,
both definitions lead to the same collection of open subsets in $\banX$,
so $\topT \eqdef \topT_{\norm{\cdot}} \eq \topT_d$,
which justifies to speak of \textit{the} topology of a normed linear space.
Clearly,
equivalent norms on $\banX$ own one and the same topology.
For example,
if $\banX$ is a finite-dimensional linear space,
then all norms on $\banX$ are equivalent and each of them
generates the so-called \textit{standard topology} on $\banX$.

\subparagraph*{} %- norm-isomorphisms
if $\banX$ and $\banY$ are normed linear spaces
and $\iota$ is a \textit{norm-preserving} linear homeomorphism
between $\banX$ and $\banY$,
that is,
$\norm{\iota x}_{\banY} \eq \norm{x}_{\banX}$ for all $x \in \banX$,
we denote this mapping a \textit{norm-isomorphism} between $\banX$ and $\banY$,
and in this case both spaces are called \textit{norm-isomorphic} to each other
(we again use the notion $\banX \cong \banY$ in this case).

\subparagraph*{} %- sequentially complete normed spaces, Banach spaces
A normed linear space $\banX$ is said to be \textit{sequentially complete}
(or briefly \textit{complete}),
if for any Cauchy sequence $\{x_n\}_{n \in \N}$ in $\banX$
there is $x \in \banX$
such that $\lim_{\,n \xto \infty} \norm{x_n-x} \eq 0$.
Equivalently,
the space $\banX$ is complete
\iff
for every absolutely convergent series
$\sum_{n=1}^{\infty} \norm{x_n} < \infty$ there is $x \in \banX$
such that $\norm{x-x_n} \xto 0$ for $n \xto \infty$.
Each complete normed space is denoted a \textit{Banach space}.
Like in the non-complete case,
a Banach space $\banX$ with norm $\norm{\,\cdot\,}$
is nothing but a {\FRECHET} space,
whose separating family of seminorms contains only this norm.

\subparagraph*{}
Notice that if the space $\banX$ is normed,
then $\banX^* \eq \BO{\banX,\F}$ is also normed,
where the norm $\norm{\phi}$ of a functional $\phi \in \banX^*$ is defined by
\[
\norm{\phi} \, \eqdef \, \sup_{x \in \banX,\,\norm{x} \leq 1} \, \abs{\phi(x)} .
\]
Since the fields $\R$ and $\C$ are complete normed spaces as well,
the {adjoint space} of a normed linear space is also complete.
This follows from  a basic theorem in functional analysis claiming
that the space $\BO{\banX,\banY}$ of continuous linear mappings
$\phi\:\banX \xto \banY$ endowed with the \textit{operator norm}
is complete \iff $\banY$ is complete.

\paragraph*{} %- norm convergence
In the following we want to introduce concepts of convergence
for sequences in a \textit{normed} linear space $\banX$.
First of all,
a sequence $\{x_n\}_{n \in \N}$ in $\banX$
is called \textit{strongly convergent}
or simply \textit{norm convergent} to $x \in \banX$,
if it converges to $x$ {\wrt} the norm topology,
that is,
if $\lim_{\,n \xto \infty} \norm{x_n-x} \eq 0$.

\subparagraph*{} %- weak topologies
But there is a further and useful topology on $\banX$,
called the \textit{weak topology} $\topT_w$.
It is the coarsest topology on $\banX$
such that all continuous linear functionals
$\phi \in \banX^*$ remain continuous.
The notion of convergence {\wrt} this topology
is known as \textit{weak convergence}.
More precisely,
a sequence $\{x_n\}_{n \in \N}$ of elements in $\banX$
is called \textit{weakly convergent} to $x \in \banX$
(notation: $x_n \rightharpoonup x$),
if $\lim_{n \xto \infty} \phi(x_n) \eq \phi(x)$ for each $\phi \in \banX^*$.

\subparagraph*{} %- properties of the weak topology
The weak topology is also Hausdorff.
But it does not contain more open sets than the norm topology,
so $\topT_w$ is weaker than $\topT_{\norm{\cdot}}$.
In other words,
norm convergence of any sequence in $\banX$ implies weak convergence.
If $\banX$ is finite-dimensional,
however,
both topologies coincide,
but this is in general not true in arbitrary normed linear spaces.

%------------------------------------------------------------------------------%
\section{Special Properties of Banach Spaces}                                  %
%------------------------------------------------------------------------------%
In this section we briefly review two fundamental properties,
a Banach space $\banX$ may or may not have.
The first topic is \textit{uniform} and \textit{strict convexity}
of real Banach spaces,
while the second one is \textit{reflexivity}.

\subparagraph*{} %- uniformly convex spaces
A \textit{uniformly convex space} is a \textit{real} normed space $\banX$
having the following property:
for every $\epsilon > 0$ there is $\delta > 0$ 
such that
for any elements $x,y$ of the closed unit ball $B_1(\0)$ of $\banX$
the validness of the inequality $\norm{(x+y)/2} > 1- \delta$
implies $\norm{x-y} < \epsilon$.
In other words,
if the center $(x+y)/2$ of the line segment $[x,y]$
between elements $x,y$ belonging to the \textit{unit sphere}
$S_1(\0) \eqdef \partial B_1(\0)$ tends to $S_1(\0)$,
then the $x$ and $y$ tend together.
This property is possibly lost
when switching to another but equivalent norm on $\banX$.

\subparagraph*{} %- strictly convex spaces
A normed linear space $\banX$ is called \textit{strictly convex},
if the conditions
$\norm{x} \leq 1$, $\norm{y} \leq 1$ and $x \neq y$
imply
$\norm{x+y}/2< 1$.
This means,
that the \textit{inner} line segment $(x,y)$
between $x,y \in S_1(\0)$ lies completely in the interior of $B_1(\0)$,
that is
$\norm{\,\alpha x + (1-\alpha)\,y\,} < 1$ for $0 < \alpha < 1$.
Each linear subspace of a uniformly convex Banach space owns this feature.
In the light of the forthcoming theorem \ref{APPROXIMATION_THEOREM},
it is an important fact
that uniformly convex spaces are also strictly convex.

\paragraph*{} %- approximation theorem
Recall that a subset $\convC$ of a linear space is called \textit{convex},
if the line segment $[x,y]$ of any pair of points $x,y \in \convC$
belongs to $\convC$ again.
Then many approximation problems may be formulated as follows:
it is to find an element $y_0$ in a convex set $\convC \subset \banX$
such that $y_0$ has minimal distance to a given element $y \notin \convC$.
The following statement,
known as the \textit{approximation theorem},
provides sufficient conditions regarding existence and uniqueness
of a solution for that problem:

\begin{theorem}\label{APPROXIMATION_THEOREM}
Let $(\banX,\norm{\,\cdot\,})$ be a real normed linear space,
let $y \in \banX$ and
$\convC \subset \banX$ be a \textit{closed} and \textit{convex} subset.
Then we have:

\begin{itemize}[leftmargin=25pt,itemsep=2pt]
\item[(a)] \textit{Uniqueness}:
if $\banX$ is \textit{strictly convex},
then there is mostly one element $y_0 \in \convC$
such that $\norm{y_0-y} \eq \inf_{x \in \convC} \norm{y-x}$.

\item[(b)] \textit{Existence}:
if $\banX$ is a \textit{uniformly convex Banach space}
then there is a unique element $y_0 \in \vecV$
such that $\norm{y_0-y} \eq \inf_{x \in \convC} \norm{y-x}$.
\end{itemize}
\end{theorem}

\paragraph*{} %- reflexive spaces
We next introduce the notion of reflexivity of Banach spaces.
It may happen
that for a given normed linear space $\banX$
the natural embedding
$\iota\: \banX \hookrightarrow \banX^{**}$ is surjective,
and in this case
it constitutes a norm-isomorphism between $\banX$ and $\banX^*$.
Many normed linear spaces considered in functional analysis
have this remarkable property,
which is indicated by the concept of \textit{reflexivity}:
the space $\banX$ is called \textit{reflexive},
if $\iota$ is onto and consequently $\banX \cong \banX^{**}$.

\paragraph*{} %- properties of reflexive spaces
We collect some propositions concerning reflexivity:

\begin{itemize}[leftmargin=35pt]
\item[1.]
Since adjoint spaces of normed spaces are throughout Banach spaces,
every reflexive space must be complete.

\item[2.]
If a Banach space $\banX$ is norm-isomorphic to a reflexive space,
then $\banX$ must be reflexive, too.

%\item[]
%All closed linear subspaces $\clU$ of a reflexive Banach space $\banX$
%and their quotient spaces $\banX / \clU$ prove to be reflexive.

\item[3.]
Dual spaces of reflexive spaces are reflexive.

\item[4.]
Reflexivity of a Banach space $\banX$ is strongly related to reflexivity
of \textit{all} its adjoint spaces,
that is,
$\banX$ is reflexive
iff
one of the spaces
$
\banX^{*}, \banX^{**}, \banX^{***}, \ldots
$
is reflexive.
On the other hand,
if $\banX$ is not reflexive,
then all these spaces prove to be pairwise unisomorphic to each other.
\end{itemize}

\subparagraph*{} %- Milman-Pettis theorem
Important examples of reflexive Banach spaces are provided
by the \textsc{Milman-Pettis} \textit{theorem}:
each \textit{uniformly convex Banach space} proves to be reflexive.

%------------------------------------------------------------------------------%
\section{Hilbert Spaces}                                                       %
%------------------------------------------------------------------------------%
Given a real or complex linear space $\hilH$,
an \textit{inner product} in $\hilH$ is defined to be a mapping
$\scalar{\,\cdot\,}{\,\cdot\,}\: \hilH \times \hilH \xto \F$,
which is

\begin{itemize}[leftmargin=35pt]
\item[(I-1)]
linear in the first argument,

\item[(I-2)]
either linear or antilinear in the second argument (conditional upon,
whether $\F$ is the real or the complex number field),
\end{itemize}
together with the additional properties

\begin{itemize}[leftmargin=35pt]
\item[(I-3)]
$\scalar{y}{x} \eq           \scalar{x}{y}$  ($\F=\R$) rsp.
$\scalar{y}{x} \eq \overline{\scalar{x}{y}}$ ($\F=\C$),
\item[(I-4)]
$\scalar{x}{x} \geq 0$ for all $x \neq 0$,
\item[(I-5)]
$\scalar{x}{x} \eq 0$ implies $x \eq 0$.
\end{itemize}

\subparagraph*{}
Property (I-3) is denoted the \textit{Hermitean symmetry}
of $\scalar{\,\cdot\,}{\,\cdot\,}$,
while (I-4) and (I-5) together mean \textit{positive definiteness}.
The presence of the inner product gives rise to define a norm in $\hilH$
according to
\begin{equation}\label{DERIVED_NORM}
\norm{x} \, \eqdef \, \scalar{x}{x} ^{1/2} \quad \quad (x \in \hilH),
\end{equation}
so each linear space furnished with an inner product is normed
in a natural manner.
The inner product and the norm are coupled by
the \textit{Cauchy-Schwarz} \textit{inequality}
\[
\abs{\scalar{x}{y}} \, \leq \,
\norm{x} \cdot \norm{y} \quad \quad (x,y \in \hilH),
\]
which implies that the mapping
$\scalar{\,\cdot\,}{\,\cdot\,}\:\hilH \times \hilH \xto \F$ is continuous.
An inner product space $\hilH$ is called a \textit{Hilbert space}
or shortly \textit{Hilbertian},
if it is complete {\wrt} the norm defined by \eqref{DERIVED_NORM}.

\subparagraph*{} %- parallelogram equation and polarization formula
The pleasant geometric properties of inner product spaces
may be attributed to the validness of the \textit{parallelogram equation}
\begin{equation}\label{PARALLELOGRAM_EQUATION}
2 \norm{x}^2 + 2 \norm{y}^2 \= \norm{x+y}^2 + \norm{x-y}^2
\qfa x,y \in \hilH.
\end{equation}
In general,
the norm $\norm{\,\cdot\,}$ on a linear space $\banX$
can be induced by an inner product
\iff \eqref{PARALLELOGRAM_EQUATION} holds.
In this case and if $\banX$ is a complex linear space,
the inner product can be recovered from the norm
by the \textit{polarization formula}
\[
\scalar{x}{y}
\=
\frac{1}{4} \,
(\,
\norm{x + y}^2 - \norm{x-y}^2 +i \norm{x+iy}^2 - i \norm{x-iy}^2
\,).
\]
If $\banX$ is a real space,
however,
the polarization formula simplifies to
\[
\scalar{x}{y} \=
\frac{1}{4} \,
(\,
\norm{x + y}^2 - \norm{x-y}^2
\,).
\]

\paragraph*{} %- properties of Hilbert spaces
In the following we want to point out
that Hilbert spaces differ from general Banach spaces
by a couple of desirable properties:

\begin{itemize}[leftmargin=15pt,itemsep=6pt]
\item[1.] \textit{Orthogonality}.
Two elements $x,y \in \hilH$ are called \textit{orthogonal} to each other,
if $\scalar{x}{y} \eq 0$ (notation: $x \bot y$).
Any subset $\orthO \subset \hilH$ is said to be \textit{orthogonal}
or an \textit{orthogonal system},
if all elements in $\orthO$ are orthogonal to each other.
Finally,
the \textit{orthogonal space} $\setA^{\bot}$
of a subset $\setA \subset \hilH$ contains exactly the elements $y \in \hilH$,
which are orthogonal to all $x \in \setA$.
Such sets basically form closed linear subspaces of $\hilH$.

\item[2.] \textit{Projection Theorem}.
Let $\setA \subset \hilH$ be a \textit{closed} linear subspace.
Then for each $x \in \hilH$ there is exactly one $x_0 \in \setA$
such that
\[
\norm{ x - x_0 } \= \inf \, \{ \, \norm {x-y}: y \in \setA \,\}.
\]
The element $x_0$ is called the \textit{projection} of $x$ onto $\setA$
and minimizes the distance between $\setA$ and $x$.
The associated linear mapping
$\pi_{\setA}\:x \in \hilH \mapsto x_0 \in \setA$
is also denoted the \textit{projection} of $\hilH$ onto $\setA$.
In general,
by a \textit{projection} is meant
any linear mapping $\pi\:\hilH \xto \hilH$ satisfying
$\pi \circ \pi \eq \pi$ and
$\scalar{\pi x}{y} \eq \scalar{x}{\pi y}$ for all $x,y \in \hilH$.

\item[3.] \textit{Decomposition Theorem}.
Using the notation of orthogonal space $\setA^{\bot}$,
the projection of $x$ onto $\setA$ is characterized as follows:
$x_0 \eq \pi_{\setA} x$ iff $x-x_0 \in \setA^{\bot}$.
Hence each $x \in \hilH$ can be uniquely represented by the formula
$x = x_0 + x^{\bot}$,
where $x_0 \in \setA$ and $x^{\bot} \in \setA^{\bot}$,
so the entire space $\hilH$ is isometrically isomorphic
to the direct sum $\setA \, \oplus \, \setA^{\bot}$
of the closed subspaces $\setA$ and $\setA^{\bot}$.
This fact is known as \textit{decomposition theorem} in Hilbert spaces.
In other words,
each closed linear subspace $\setA$ of $\hilH$
is \textit{topologically complementable},
that is,
to each closed subspace $\setA$
there is a second closed subspace $\setB \subset \hilH$
such that $\setA \cap \setB \eq \{\0\}$ and $[\setA \cup \setB] \eq \hilH$.
Of course,
a closed subspace $\A$ is topologically complemented
by its orthogonal space. %$\setA^{\bot}$.

\item[4.] \textit{Existence of Orthonormal Bases}.
An orthogonal system $\orthO \subset \hilH$ is denoted
an \textit{orthonormal system},
if $\scalar{x}{x} \eq 1$ for all $x \in \orthO$.
However,
if $\orthO$ contains only \textit{finitely} many elements,
then the span $[\orthO]$ forms a \textit{closed} subspace in $\hilH$,
and the projection $\pi_{[\orthO]}$ of any element $x$ onto $[\orthO]$
can be expressed by the formula
\[
\pi_{[\orthO]} \, x \= \sum_{i=1}^n \, \scalar{x}{e_i} \, e_i.
\]
However,
if $\orthO$ consists of arbitrarily many elements,
the inequality $\scalar{x}{e_{\alpha}} \neq 0$ can only be valid
for mostly countable many elements $e_{\alpha} \in \orthO$,
which in conclusion enables to add the coefficients $\scalar{x}{e_{\alpha}}$
together about all elements of $\orthO$.

\subparagraph*{} %- orthonormal bases
An orthonormal system $\orthO$ is called
an \textit{orthonormal base} in $\hilH$,
if its span $[ \orthO ]$ lies dense in $\hilH$.
In this case we may represent each $x \in \hilH$ according to
\begin{equation}\label{FOURIER_COEFFICIENTS}
x \= \sum_{e_{\alpha} \in \orthO} \, \scalar{x}{e_{\alpha}} \, e_{\alpha},
\end{equation}
which is primarily a formal sum,
but consists of mostly countable many summands.
The inner products $\scalar{x}{e_{\alpha}}$ are called
the \textit{Fourier coefficients} of $x$ {\wrt} $\orthO$.
The right hand side of \eqref{FOURIER_COEFFICIENTS}
consists of \textit{mostly countable many} non-zero summands,
and this series is invariant
by realignments of the elements $e_{\alpha} \in \orthO$.

\subparagraph*{}
There is of course the need to determine,
whether a given orthonormal system $\orthO$ forms an orthonormal base.
In fact,
it can be shown that
the following statements are equivalent:

\begin{itemize}[leftmargin=20pt,itemsep=3pt]
\item[(a)]
$\orthO$ is an orthonormal base in $\hilH$.

\item[(b)]
$\orthO$ is \textit{complete},
that is,
$\scalar{x}{e} \eq 0$ for all $e \in \orthO$ implies $x \eq \0$.

\item[(c)]
$\sum_{e_{\alpha} \in \,\orthO} \;
\abs{ \scalar{x}{e_{\alpha}} } ^2 \eq \norm{x}^2$ for all $x \in \hilH$
(\textit{Parseval}\textit{'s equation}).
\end{itemize}

\paragraph*{} %- existence of orthonormal bases
It is also nearby to ask for the existence of orthonormal bases.
First of all,
the collection of all orthonormal systems in a Hilbert space $\hilH$
forms a partial ordered set,
and under the assumption that \textsc{Zorn}{'s lemma}\,\footnote{\,
\,The \textit{axion of choice} claims that
if $\cal{F}$ is a disjoint family $F_{\alpha}$ ($\alpha \in \J$)
of non-empty sets,
then there is a set $A$ that shares exactly one element
with each other set $F_{\alpha} \in \cal{F}$.}
is valid,
it contains a \textit{maximal} element $\orthO$
that cannot be increased,
hence each Hilbert space $\hilH \neq \{\0\}$ owns an orthonormal base.
%\subparagraph*{} %- Hilbert space dimension
Furthermore,
if $\orthO_1$ and $\orthO_2$ are orthonormal bases in $\hilH$,
then there is a bijective mapping between them,
thus they own identical cardinal numbers.
This legitimates us to introduce the notation
of \textit{Hilbert space dimension} $\abs{\hilH} $
being the cardinal number of any orthonormal base in $\hilH$.

\subparagraph*{} %- separable Hilbert spaces
A Hilbert space $\hilH$ is called \textit{separable}
if there is a mostly countable subset lying dense in $\hilH$.
It turns out that $\hilH$ is separable
\iff
$\abs{ \hilH} \leq \abs{\N} $.
Indeed,
each orthonormal base in $\hilH$ forms a mostly countable set in this case,
and $\hilH$ proves to be norm-isomorphic to the Hilbert space $\ell^2$
of square summarizable sequences $f\:\N \xto \F$.
\end{itemize}

\subsubsection*{} %- hierarchy of topological linear spaces
The following figure shows the inheritance tree
of topological and normed linear spaces,
where Banach and Hilbert spaces are positioned at the very end
of this hierarchy:

\begin{center}
\scalebox{0.69}{

\begin{tikzpicture}[node distance=4.0cm,auto]

\node[align=center,rectangle,draw](TLS) at (0,-1.0)
{\;topological linear space\; \\ $(\topV,\topT_V)$};

\node[align=center,rectangle,draw](LCTLS) at (-5,-3.5)
{locally convex \\ \;topological linear space\; \\ $(\topV,\{ p_{\alpha} \}_{\alpha \in \cal{I}})$};

\node[align=center,rectangle,draw](SCTLS) at (5,-3.5)
{\textbf{sequentially complete}\\ \; \textbf{topological linear space} \; \\ $(V,\topT_V)$};

\node[align=center,rectangle,draw](MS) at (-3,-5.9)
{metric space \\ $(\metrS,\topT_d)$};

\node[align=center,thick,rectangle,draw](CMS) at (1,-5.9)
{\textbf{complete} \\ \; \textbf{metric Space} \; \\ $(\metrS,\topT_d)$};

\node[align=center,thick,rectangle,draw](F-S) at (5,-5.9)
{\; \textbf{F-space}\; \\ $(F,\topT_V=\topT_d)$};

\node[align=center,rectangle,draw](PFS) at (-5,-8.0)
{\;pre-{\FRECHET} space\; \\ $(F,\{p_n\}_{n \in \N})$};

\node[align=center,thick,rectangle,draw](FS) at (5,-8.0)
{\;\textbf{{\FRECHET} space}\; \\ $(F,\{p_n\}_{n \in \N})$};

\node[align=center,rectangle,draw](NS) at (-5,-10.5)
{\;{normed linear space}\; \\$(\banX,\norm{\cdot})$};

\node[align=center,thick,rectangle,draw](BS) at (5,-10.5)
{\;\textbf{Banach space}\; \\$(\banX,\norm{\cdot})$};

\node[align=center,thick,rectangle,draw](RBS) at (5,-12.0)
{\;\textbf{reflexive Banach space}\; \\$(\banX,\norm{\cdot})$};

\node[align=center,thick,rectangle,draw](UCS) at (5,-13.5)
{\;\textbf{uniformly convex space}\; \\$(\banX,\norm{\cdot})$};

\node[align=center,rectangle,draw](IPS) at (-5,-15)
{\;{inner product space}\; \\$(\vecH,\scalar{\cdot}{\cdot})$};

\node[align=center,thick,rectangle,draw](HS) at (5,-15)
{\;\textbf{Hilbert space}\; \\$(\hilH,\scalar{\cdot}{\cdot})$};

\draw[->,>={Stealth}, thick] (MS) to (CMS);
\draw[->,>={Stealth}, thick] (MS) to (PFS);
\draw[->,>={Stealth}, thick] (CMS) to (F-S);
\draw[->,>={Stealth}, thick] (TLS) to (SCTLS);
\draw[->,>={Stealth}, thick] (TLS) to (LCTLS);
\draw[->,>={Stealth}, thick] (SCTLS) to (F-S);
\draw[->,>={Stealth}, thick] (F-S) to (FS);
\draw[->,>={Stealth}, thick] (LCTLS) to (PFS);
\draw[->,>={Stealth}, thick] (PFS) to (NS);
\draw[->,>={Stealth}, thick] (NS) to (IPS);
\draw[->,>={Stealth}, thick] (PFS) to (FS);
\draw[->,>={Stealth}, thick] (NS) to (BS);
\draw[->,>={Stealth}, thick] (IPS) to (HS);
\draw[->,>={Stealth}, thick] (FS) to (BS);
\draw[->,>={Stealth}, thick] (BS) to (RBS);
\draw[->,>={Stealth}, thick] (RBS) to (UCS);
\draw[->,>={Stealth}, thick] (UCS) to (HS);

\node[align=center](I) at (0,0.4) {\textit{metric and topological linear spaces}};
\node[align=center](II) at (   0,-16.4) {\textit{normed linear spaces}};
\draw[dashed,thick]        (-7.8,  0.0) rectangle (8.0,-9.0);
\draw[dashed,thick]        (-7.8, -9.5) rectangle (8.0,-16.0);
\end{tikzpicture}
}
\end{center}

\begin{center}
\footnotesize{Figure \thesection.1:
              The hierarchy of metric and topological linear spaces}
\end{center}

%------------------------------------------------------------------------------%
\section{Bochner-Sobolev Spaces}                                               %
%------------------------------------------------------------------------------%
This section is preliminary in order to outline the variational approach
to evolution equations in the Hilbert space setting.

\subparagraph*{} %- Bochner spaces
The \textit{Bochner spaces} $L^p(\Omega,\banX)$ are nothing but
Lebesgue spaces for vector-valued functions $f\:\Omega \xto \banX$,
where $\banX$ is a separable Banach space.
To define them,
we first have to introduce a concept of measurability and integration
for vector-valued functions,
that forms a generalization of the ordinary Lebesgue integral.

\subparagraph*{} %- measurable functions
A function $f\:\Omega \xto \banX$ is called \textit{weakly measurable},
if $\phi \circ f\:\Omega \xto \F$ is measurable for every $\phi \in \banX^*$.
Then a weakly measurable function $f\:\Omega \xto \banX$
is called \textit{Bochner-integrable},
if
\[
\int_{\Omega} \norm{f} \,d\mu \, < \, \infty.
\]
For any $p \geq 1$
we denote $L^p(\Omega,\banX)$ the space of functions $f\:\Omega \xto \banX$
having Bochner-integrable powers $\abs{f} ^{p}$.
The elements of this space have also to be understood
as equivalence classes of functions
being identical almost everywhere in $\Omega$.
Endowed with the norm
\[
f \mapsto \norm{f}_p \eqdef (\int_{\Omega} \, \norm{f} ^p \, d\mu)^{1/p},
\]
the Bochner spaces $L^p(\Omega,\banX)$ are complete
and thus constitute Banach spaces.

\subparagraph*{} %- weak and strong measurability
It should be mentioned that
if we suppose the Banach space $\banX$ to be separable,
then a function $f\:\Omega \xto \banX$ is weakly measurable
if and only if $f$ is strongly (or Pettis-) measurable.
The latter notion of measurability uses approximations of $f$ by sequences
$\{e_n\}_{n \in \N}$ of elementary functions $e_n\:\Omega \xto \banX$,
and the Bochner integral of $f$ is obtained by taking the limit of the
Bochner integrals of the functions $e_n$,
that is,
\[
\int_{\Omega} f \,d\mu \eqdef \lim_{n \xto \infty} \int_{\Omega} e_n \, d\mu.
\]
This definition proves to be independent of the specific choices
of the approximating sequences for $f$.

\subparagraph*{} %- specific choice of Omega and X for Bochner spaces
With respect to applications in partial differential equations,
Bochner spaces are of interest on the one hand 
for the choices $\Omega \eq \G$ and $\banX \eq \F \in \{\R,\C\}$,
where $\G$ is a domain in one of the Euclidean spaces $\Rd$
and $L^p(\G,\F)$ is nothing but the Lebesgue space $\Lp$.
On the other hand,
in the weak solution theory of evolution equations
we chose $\Omega$ to be a finite or infinite interval $\J \subset \R$,
and the Banach space $\banX$ may be any complete function space $\mathbb{F}$
consisting of differentiable mappings $f\:\G \xto \F$,
so we are dealing with the Bochner spaces $L^p(\J,\mathbb{F})$ in this case.

\paragraph*{} %- Bochner-Sobolev spaces 
We finally introduce so-called \textit{Bochner-Sobolev spaces},
which differ by the property
that their members are at the same time
Bochner-integrable and \textit{weakly differentiable}.
Let $\hilH$ be a separable Hilbert space and $L^2(a,b;\hilH)$ be
the Bochner space consisting of functions $f\:(a,b) \xto \hilH$,
that is,
we have
\[
\int_a^b \norm{f(t)}^2 \,dt \, < \, \infty
\]
for $f \in L^2(a,b,;\hilH)$.
We endow this space with the inner product
\[
\scalar{f}{g} \eqdef \int_a^b \scalar{f(t)}{g(t)} \,dt,
\quad \quad
f,g \in L^2(a,b;\hilH)
\]
such that the norm is given as usual by
$\norm{f}_2 \eq \scalar{f}{f}^{1/2}$.
Thus $L^2(a,b;\hilH)$ is Hilbertian,
and if $(a,b)$ is a \textit{finite} interval,
the inclusion $L^2(a,b;\hilH) \subset L^1(a,b;\hilH)$ holds.
We shall denote $L^2(a,b;\hilH)$ a \textit{Hilbert-Bochner space}.

\subparagraph*{} %- weak derivatives of vector-valued functions
For $f \in L^2(a,b;\hilH)$,
a further function $g\:(a,b) \xto \hilH$ is said to be
the \textit{weak derivative} of $f$,
if
\[
\int_a^b g \phi \,dx \= - \int_a^b f \phi' \,dx
\]
holds for each $\phi \in C_c^{\infty}(a,b,\banX)$,
that is,
$\phi$ is an indefinitely differentiable function with compact support
in $(a,b)$.
If such a function $g$ exists,
we denote $f' := g$ in this case,
which should not be confused with the classical derivative of $f$.
But notice that if the weak derivative of $f$ is present,
then $f$ is also classical differentiable.

\subparagraph*{} %- Hilbert-Bochner-Sobolev space
Having the notion of weak derivative available,
we may introduce the \textit{Hilbert-Bochner-Sobolev space}
\[
H^1(a,b;\hilH) :=
W^{1,2}(a,b,\hilH) :=
\{f \in L^2(a,b;\hilH): f' \in L^2(a,b;\hilH) \},
\]
which is a Hilbert space as well with inner product given by
\[
\scalar{f}{g}_1 \eqdef
\int_a^b \scalar{f}{g} \, dt
\, + \,
\int_a^b \scalar{f'}{g'} \, dt\,,
\quad f,g \in H^1(a,b;\hilH).
\]

\subparagraph*{} %- mixed Bochner-Sobolev spaces
Finally,
let $\hilH_1, \hilH_2$ be \textit{real} Hilbert spaces,
where $\hilH_1$ is continuously embedded into $\hilH_2$.
Then there is also a continuous embedding
\[
H^1(a,b;\hilH_1) \, \overset{\kappa}{\hookrightarrow} \, H^1(a,b;\hilH_2).
\]
Consequently,
given a Hilbert space triple 
$
\vecV \, {\overset{\iota_1}{\hookrightarrow}} \,
\hilH \, {\overset{\iota_2}{\hookrightarrow}} \vecV^*,
$
we may regard the space $H^1(a,b;\vecV)$
as a linear subspace of $H^1(a,b;\vecV^*)$,
which enables us to define the \textit{mixed Bochner-Sobolev space}
\[
H^1(a,b;\vecV,\vecV^*) \eqdef
\{ f \in L^2(a,b;\vecV): f' \in L^2(a,b;\vecV^*) \}
\]
or equivalently
\[
H^1(a,b;\vecV,\vecV^*) \eqdef L^2(a,b;\vecV) \, \cap H^1(a,b;\vecV^*).
\]
Then the mapping
\[
(f,g) \mapsto \scalar{f}{g}_1 \eqdef
\int_a^b \, \scalar{f(t)}{g(t)}_{\vecV} \, dt
\, + \,
\int_a^b \, \scalar{f'(t)}{g'(t)}_{V^*} \, dt
\]
fulfills the properties on an inner product
and makes $H^1(a,b;\vecV,\vecV^*)$ into a Hilbert space
having the following properties:

\begin{itemize}[leftmargin=7mm]
\item[(a)]
Each $f \in H^1(a,b;\vecV,\vecV^*)$
is $\lambda$-nearly everywhere identical
to a continuous function $f_c\:[a,b] \xto \hilH$,
and by identifying the functions $f$ and $f_c$
the space $H^1(a,b;\vecV,\vecV^*)$ can continuously be embedded
into the space $C([a,b],\hilH)$.

\item[(b)]
The space $C^{\infty}([a,b],\vecV)$ is dense in $H^1(a,b;\vecV,\vecV^*)$.

\item[(c)]
The following product rule is valid for each $t \in (a,b)$:
\[
\frac{d}{dt} \norm{u(t)} ^2
\=
2 \re \, \scalar{u'(t)}{u(t)}_1
\fa u \in H^1(a,b;\vecV,\vecV^*).
\]
\end{itemize}

%------------------------------------------------------------------------------%
\section{Applications}                                                         %
%------------------------------------------------------------------------------%
In this section let $\G$ be a \textit{domain} in $\Rd$,
that is,
an open and connected set within an Euclidean space $\Rd$,
which may be bounded, unbounded or even the space $\Rd$ itself.
The application part of this section consists mainly of an enumeration
of function spaces having meaning in PDE theory and their adjoints,
which is followed by some examples of orthonormal bases
in Hilbert space $L^2(\J)$,
where $\J \subset \R$ is an interval.

\subparagraph*{} %- multi-indices
Preliminary,
we will work in this section with \textit{multi-indices},
which are nothing but $d$-tuples
$\alpha \eq (\alpha_1,\ldots,\alpha_d)$
with values $\alpha_i \in \N_0$
and modulus $\abs{\alpha} \eqdef \alpha_1 + \ldots + \alpha_d$.
They are mostly used to shorten multiple partial derivatives
$
\PA f \eqdef \partial_1^{\alpha_1}
         \partial_2^{\alpha_2} \ldots
         \partial_d^{\alpha_d} f
$
of a sufficiently often differentiable and complex-valued function $f$,
but also to denote the multiple powers
$\xi^{\alpha} \eqdef \xi_1^{\alpha_1}
\cdot \xi_2^{\alpha_2} \cdot \ldots \cdot \xi_d^{\alpha_d}$
of a vector $\,\xi \eq (\xi_1,\xi_2,\ldots,\xi_d) \in \C^d$.

\begin{itemize}[leftmargin=15pt,itemsep=6pt]
\item[1.] \textit{Spaces of continuous functions}.
The space $C(\clG)$ consists of functions $f\:\G \xto \F$,
which are restrictions of continuous functions $\tilde{f}\:\G' \xto \F$,
where $\G'$ is an open superset of $\G$
such that $\clG \subset \G'$
(notation: $\G \Subset \G'$).
It is a well-known fact that
$C(\clG)$ is complete {\wrt} the \textit{supremum norm}
\[
\norm{f}_{\infty} \, \eqdef \, \sup_{x \in \G} \abs{f(x)} .
\]
This is true as well for the linear subspace $C_0(\clG) \subset C(\clG)$,
consisting of functions that are vanishing at infinity,
that is,
for each $\epsilon > 0$
there is a compact set $C_{f,\epsilon} \subset \G$
such that
$\abs{f(x)} < \epsilon$
for all points $x \in \G \setminus C_{f,\epsilon}$.
%The space $C(\clG)$ is not reflexive.

\item[2.] \textit{Spaces of continuously differentiable functions}.
The linear subspace $C^k(\clG) \subset C(\clG)$
consists of those continuously differentiable functions $f\:\G \xto \F$,
whose partial derivatives $\partial^{\alpha} f$ are bounded on $\clG$
for every multi-index $\alpha$ with $\abs{\alpha} \leq k$.
It is complete if endowed with the norm
\[
\norm{f}_k
\, \eqdef \,
\sum_{\abs{\alpha} \leq k} \,
\sup_{x \in \clG } \; \abs{\partial^{\alpha} f(x)} .
\]
The spaces $C^k(\clG)$ are not reflexive.

\item[3.] \textit{{\HOELDER} spaces}.
It has been shown
that the spaces $C^k(\clG)$ are not very convenient
for the investigation of partial differential equations,
but the so-called \textit{{\HOELDER} spaces} $C^{k,\gamma}(\clG)$ are so.
Preliminary,
a function $f\:\G \xto \F$ is called \textit{{\HOELDER}-continuous}
{\wrt} the parameter $\gamma \in (0,1]$,
if its \textit{{\HOELDER}-norm} is finite,
that is,
\[
\abs{f} _{\,\gamma} \, \eqdef \,
\sup_{x,y \, \in \, \G, \, x \neq y} \,
\frac{\abs{f(x) - f(y)} } {\abs{x-y} ^{\gamma}} \, < \, \infty.
\]
Let $C^{k,\gamma}(\clG)$ be the linear space
of functions $f \in C^k(\clG)$,
whose partial derivatives $\PA f$
are {\HOELDER}-continuous for all $\abs{\alpha} \eq k$.
Then the \textit{{\HOELDER} norm} $\abs{f} _{k,\gamma}$
of a function $f \in C^{k,\gamma}(\clG)$ is defined
by the sum of the norm $\norm{f}_k$
and the {\HOELDER} norms of the partial derivatives $\PA f$
for $\abs{\alpha} \eq k$:
\[
\abs{f} _{k,\gamma}
\, \eqdef \,
\norm{f}_k
\, + \,
\sum_{\abs{\alpha} \eq k} \abs{\,\PA f} _\gamma.
\]
The space $C^{k,\gamma}(\clG)$ is complete {\wrt} this norm,
but it is not a reflexive one.
For $k=0$ and $\gamma \eq 1$ we obtain the space
$C^{0,1}(\clG)$ of \textit{Lipschitz-continuous} functions.

\subparagraph*{}
If $\G$ is a bounded domain, then we have for $0 < \alpha < \beta \leq 1$
\[
C^1(\clG) \subsetneq
C^{0,\beta}(\clG) \subsetneq
C^{0,\alpha}(\clG) \subsetneq
C(\clG).
\]

\item[4.] \textit{Spaces of smooth functions}.
Next we consider three function spaces
consisting of indefinitely differentiable functions $f\:\G \xto \F$
that frequently occure in applied analysis
and constitute {\FRECHET} spaces respectively
complete locally convex linear spaces.

\begin{itemize}[leftmargin=20pt]
\item[(a)]
The space $\SF$ of functions $f\:\G \xto \F$ being indefinitely differentiable.
Unfortunately,
there is no norm that makes $\SF$ into a complete normed space,
but after all it is possible to define a countable family $\famF$
of seminorms to make $\SF$ into a {\FRECHET} space.
For example,
$\famF$ can be given by the mappings
\[
p_{N,n} \, \eqdef \,
\max_{x \in \compK_n} \;
\{ \, \abs{\, (\PA \phi)(x) \,} : \abs{\alpha} \leq N \}
\quad \quad (N,n \in \N)\,,
\]
where $(\compK_n)$ is a sequence of compact subsets of $\G$
such that
\[
\compK_1 \subset \compK_2 \ldots \subset \G \, \mbox{ for } \, n \in \N,
\quad
\bigcup_{n \in \N} \, \compK_n \= \G,
\]
and each $\compK_n$ lies in the interior of $\compK_{n+1}$
(such a sequence is called a \textit{compact exhaustion} of $\G$).
The space $\SF$ is often denoted $\E(\G)$,
when considered in the locally convex topology induced by the totality
of seminorms $p_{N,n}$.
It is complete {\wrt} this topology,
so it forms a {\FRECHET} space.

\item[(b)]
The subspace $\TF \subset \SF$ of smooth functions
having \textit{compact} support
\[
\supp{\,\phi} \, \eqdef \, \overline{\set{x \in \G}{\phi(x) \neq 0}}\,,
\]
usually denoted the \textit{space of test functions}.
This function space plays an outstanding role in analysis,
since it lies dense in a multiplicity of other function spaces.
Unfortunately,
there is neither a countable family of seminorms
making $\TF$ into a {\FRECHET} space nor is $\TF$ completely metrizable.

\subparagraph*{}
But it is still possible to define a locally convex topology $\topT$ on $\TF$
such that $\TF$ is complete in the sense of a topological linear space.
Instead of saying how the open sets of this topology $\topT$ look like,
let us define $\topT$ by means of a suitable notion of convergence,
known as $\D$-convergence in $\TF$:\,
a sequence $(\phi_n)$ in $\TF$ is convergent to $\phi \in \TF$,
if there is a compact set $\compK \subset \G$
such that
(i)  $\supp{\phi_n} \subset \compK$ for all $n \in \N$ and
(ii) for each multi-index $\alpha$ the sequence $(\PA \phi_n)$
     converges to $\PA \phi$ \textit{uniformly} on $\compK$.

The space $\TF$ equipped with the topology
induced by $\D$-convergence is usually denoted $\D(\G)$.
Notice
that we have $\D(\G) \hookrightarrow \E(\G)$,
that is,
the space $\D(\G)$ is continuously embedded into the space $\E(\G)$
{\wrt} the locally convex topologies of these spaces.

\subparagraph*{}
We have $\D(\G) \hookrightarrow \E(\G)$,
that is,
the space $\D(\G)$ is continuously embedded into the space $\E(\G)$
{\wrt} their locally convex topologies.

\item[(c)]
The \textit{Schwartz space} $\SRd$,
which is the space of all smooth functions $\phi \in C^{\infty}(\Rd)$
having the additional property
\[
\abs{ \PA \phi(x) \cdot x^{\beta}} \xto 0 \,
\for \norm{x}_d \xto \infty
\]
and arbitrary multi-indices $\alpha,\beta$,
where $\norm{x}_d$ is the Euclidean norm of $x \in \Rd$.
The elements of $\SRd$ are called \textit{rapidly decreasing}
or briefly \textit{Schwartz functions}.
Hence a Schwartz function and all its derivatives vanish at infinity
faster then the reciprocal of any polynomial.
The most famous example of a Schwartz function is the Hermite function
\[
\cal{H}_0(x) \, \eqdef \, e^{- \frac{\norm{x}^2}{2}} \quad \quad (x \in \Rd).
\]

\subparagraph*{}
Rather similar to the space $\SF$,
it is possible to find for the Schwartz space $\SRd$
a family of seminorms as follows:
define for $\phi \in C^{\infty}(\Rd)$ and multi-indices $\alpha, \beta$
\begin{equation}\label{I_SRN_FAMILY_OF_SEMINORMS}
p_{\alpha,\beta}(\phi) \, \eqdef \,
\sup_{x \in \Rd} \abs{ \, x^{\alpha} (D^{\beta} \phi) (x) \, } .
\end{equation}

\subparagraph*{}
Then $\phi$ is a Schwartz function
iff $p_{\alpha,\beta}(\phi) < \infty$
holds for all indices $\alpha, \beta$.
The collection of mappings $p_{\alpha,\beta}\:\SRd \xto \R$
defines a countable family of seminorms,
which gives rise to a locally convex topology $\topT_{\mathscr{S}}$ on $\SRd$.
Moreover,
since the Schwartz space is complete {\wrt} $\topT_{\mathscr{S}}$,
it actually forms a \textit{{\FRECHET} space}.

\subparagraph*{}
Comparing the Schwartz space $\SRd$ with the function spaces
$\D(\Rd)$ and $\E(\Rd)$,
it is easy to see that there is the chain of inclusions
\[
\D(\Rd) \, \hookrightarrow \, \SRd \, \hookrightarrow \, \cal{E}(\Rd).
\]
Moreover,
by comparing the locally convex topologies of these spaces,
it turns out that these inclusions are continuous.
\end{itemize}

\item[5.] \textit{Lebesgue spaces}.
To introduce the spaces $L^p(\G)$ for $p \in [1,\infty)$,
let $\cal{L}^p(\G)$ be the $\F$-linear space of measurable functions
$f\:\G \xto \F$
such that $\abs{f} ^p$ is Lebesgue-integrable.

\subparagraph*{}
We define for $f,g \in \cal{L}^p(\G)$ the equivalence relation
$f \sim g$ according to $f=g$ almost everywhere.
Then the elements of the space $L^p(\G)$
are the equivalence classes in $\cal{L}^p(\G)$ {\wrt} $\sim$,
that is,
$L^p(\G) \eqdef \cal{L}^p(\G) / \sim$.
The \textit{Lebesgue norm} in $\Lp$ is given by
\[
\norm{f}_p
\, \eqdef \,
\big( \int_{\,\G} \abs{f} ^p \, d\mu \,\big) \, ^{\frac{1}{p}}.
\]
Finally,
the \textsc{Fischer-Riesz} \textit{theorem} claims
that the spaces $L^p(\G)$ are complete {\wrt} the Lebesgue norm.
Moreover,
they are reflexive for $p >1$,
but $L^1(\G)$ is not so.

\subparagraph*{} %- Lebesgue space L2
The case $p \eq 2$ gives rise to a special attention:
here it is possible to derive the norm of the space $L^2(\G)$
from the inner product
\[
\scalar{f}{g} \, \eqdef \, \int_{\,\G} f \, \overline{g} \; d\mu,
\]
thus $\norm{f}_2 \eq \scalar{f}{f} ^{1/2}$
and $L^2(\G)$ becomes into a Hilbert space.

\subparagraph*{}
For $p=\infty$ and any measurable function $f\:\G \xto \F$ we define the mapping
\begin{equation}\label{SUPREMUM_NORM_II}
\norm{f}_{\infty} \, \eqdef \, \inf
\{\,a \geq 0: \abs{f} \leq a \,\, \mu \mbox{-almost everywhere}\,\},
\end{equation}
taking values in the set $\R^+ \cup \{\infty\}$.
Let $\cal{L}^{\infty}(\G)$ be the linear space
of all measurable functions $f\:\G \xto \F$
with $\norm{f}_{\infty} < \infty$
and $\Li$ be the linear space of all equivalence classes
in $\cal{L}^{\infty}(\G)$,
where $f \sim g$ means $f=g$ $\mu$-almost everywhere.
Members of these spaces are called \textit{essentially bounded}.
Similar to the case $p < \infty$,
$\Li$ is complete {\wrt} \eqref{SUPREMUM_NORM_II}.
But notice that this space is not reflexive.

\item[6.] \textit{Sobolev spaces}.
A function $f\:\G \xto \C$ is said to belong to the space $\Lc$,
if $f$ is Lebesgue-integrable on each compact subset of $\G$.
This space is useful to define the concept of \textit{weak differentiability}.
A function $f \in \Lc$ is said to be \textit{weakly differentiable}
{\wrt} a given multi-index $\alpha$,
if there is $g \in \Lc$
such that 
\[
\int_{\,\G} g(x) \, \phi(x) \, dx
=
(-1)^{ \abs{\alpha} } \,
\int_{\,\G} f(x) \, \PA \phi(x) \, dx
\for \phi \in \TF.
\]
We denote $\PA f \eqdef g$ in this case,
and weak differentiability arises as a linear mapping
$\PA\:D_{\alpha} \subset \Lc \xto \Lc$,
where each $D_{\alpha}$ is a linear subspace of $\Lc$.
Notice that
the concept of weak derivative is compliant
with that of the classical derivative:
if $f$ has classical derivative $\PA f$,
then $f$ is also weakly differentiable,
and the weak derivative is equal to $\PA f$.

\subparagraph*{}
The next function spaces under consideration are those,
whose elements are $p$-integrable and weakly differentiable up to order $k$
at the same time.
More precisely,
for any integer $k$ and real $p \geq 1$,
the \textit{Sobolev space} $\SobW$ is the set
of functions $f$ in the Lebesgue space $\Lp$,
whose weak derivatives $\PA f$ exist and also belong to $\Lp$
for all multi-indices $\alpha$ with $ \abs{\alpha} \leq k$,
briefly written
\[
\SobW
\, \eqdef \,
\{\,f \in \Lp: \PA f \in \Lp, \; \abs{\alpha} \leq k \, \}.
\]
Sobolev spaces are of course linear spaces,
and especially for $k=0$
the spaces $W^{0,p}(\G)$ are nothing but the Lebesgue spaces $\Lp$ themselves.

\subparagraph*{} %- Sobolev norms
For functions $f \in \SobW$ let the mapping $f \mapsto \abs{f} _{k,p}$
be defined according to
\begin{equation}\label{SOBOLEV_NORM}
\abs{f} _{k,p} \, \eqdef \,
\big(
\sum_{\abs{\alpha } \leq k}
\int_{\,\G} \, \abs{\PA f} ^p \, dx
\big) ^{1/p}.
\end{equation}
It fulfills all the properties of a norm,
thus the space $\SobW$ is actually a normed space,
and the mapping defined by \eqref{SOBOLEV_NORM}
is usually denoted a \textit{Sobolev norm}.
With respect to the question of completeness,
it is not hard to show that
the spaces $\SobW$ are complete
and therefore constitute Banach spaces as well.

\subparagraph*{} %- alternative definition of Sobolev space
There is a second way to define Sobolev space.
Consider the subspace
\[
C_{k,p}^{\infty}(\G) \, \eqdef \, \{\, f \in \SF: \abs{f} _{k,p} < \infty \,\},
\]
and since this space is \textit{not} complete
{\wrt} the norm $\abs{\,\cdot\,} _{k,p}$,
we set
\[
\cal{W}^{k,p}(\G) \, \eqdef \, \overline{C_{k,p}^{\infty}(\G)}.
\]
The space $\cal{W}^{k,p}(\G)$ is defined by a purely topological operation,
that is,
by taking the closure of a smaller space.
Thus for elements in this space there is no instrinsic notion of derivative,
and one introduces the notion of \textit{strong derivative} for
$u \in \cal{W}^{k,p}(\G)$ as follows: 
given a multi-index $\alpha$ ($\abs{\alpha} \leq k$),
a further function $\PA u$ also belonging to $\cal{W}^{k,p}(\G)$
is said to be the $\alpha$-strong derivative of $u$,
if there is a sequence $\{\phi_n\} \in C_{k,p}^{\infty}(\G)$
converging to $u$ such that
\[
\abs{\PA u - \PA \phi_n)} _{k,p} \xto 0 \quad \quad (n \xto \infty).
\]
Sobolev spaces are reflexive for $1 < p < \infty$,
but not for $p \eq 1$.

\paragraph*{} %- Meyers-Serrin theorem
The \textsc{Meyers-Serrin} \textit{theorem} claims
$\cal{W}^{k,p}(\G) \cong \SobW$ in the sense of Banach space isomorphisms.
Moreover,
the notions of strong and weak derivative are equivalent
when identifying functions in $\cal{W}^{k,p}(\G)$ and $\SobW$.

\paragraph*{} %- embedding theorems for Sobolev spaces
One of the most important properties of Sobolev spaces is given by the fact
that under certain hypotheses they can continuously be embedded
into further spaces and those of type $C^k(\clG)$.
We define for fixed $1 \leq p < \infty$
the \textit{Sobolev numbers} $\theta_{\,k,d,p}$ by
\[
\theta_{\,k,d,p} \eqdef k - d/p
\quad \quad (k,d \in \N, p \geq 1),
\]
which play a crucial role in \textsc{Sobolev}\textit{'s embedding theorem}:

\begin{theorem}\label{SOBOLEV_EMBEDDING_THEOREM_I}
Let $\G \subset \Rd$ be a domain having Lipschitz-continuous boundary.
If $k,m \in \N$ are chosen such that $\theta_{\,k,d,p} > m$,
then there are continuous embeddings
$\SobW \hookrightarrow C^m(\clG)$ and
$\SobW \hookrightarrow C^{m,\gamma}(\clG)$ for $0 < \gamma < 1$.
\end{theorem}

Hence each equivalent class $f \in \SobW$ has a representative in these spaces,
supposed the Sobolev number is sufficiently large.

\paragraph*{} %- embeddings of Sobolev spaces into other Sobolev spaces
In case of $\theta_{\,k,d,p} \leq 0$,
however,
one cannot embed $\SobW$ into $C(\clG)$.
But it is at least possible to find an embedding of $\SobW$
into some lower-dimensional Lebesgue spaces $L^q(\G)$:

\begin{theorem}\label{SOBOLEV_EMBEDDING_THEOREM_II}
For $\G \subset \Rd$,
$k \in \N$ and $p,q \geq 1$ such that
\[
1 < q < \frac{dp}{d-kp},
\]
there are continuous embeddings $\SobW \hookrightarrow L^q(\G)$.
\end{theorem}

\subparagraph*{} %- embeddings for fractional Sobolev spaces
Finally,
there are also embeddings of Sobolev spaces into further Sobolev spaces,
where the parameter $p$ is fixed:

\begin{theorem}\label{SOBOLEV_EMBEDDING_THEOREM_III}
If $\G \subset \Rd$ has Lipschitz-continuous boundary,
$s,s' \in \N$ and $1 < s < s' < \infty$,
then $W^{s',p}(\G) \hookrightarrow W^{s,p}(\G)$.
\end{theorem}

\item[7.] \textit{Hilbert-Sobolev spaces}.
Let us consider in the family $\SobW$ of Sobolev spaces
the case $p{=}2$ more specifically and denote the space
$H^k(\G) \eqdef W^{k,2}(\G)$ for $k \in \N_0$,
which is often called a \textit{Hilbert-Sobolev space}.

\subparagraph*{} %- inner product and norm of Hilbert-Sobolev spaces
Each space $H^k(\G)$ can be endowed with the inner product
\[
\scalar{f}{g}_k \, \eqdef \,
\sum_{\abs{\alpha} \leq k} \int_{\,\G} \;
\PA f \; \PA \overline{g} \, dx\,,
\quad \quad f,g \in H^k(\G)
\]
and with norm
$\abs{f} _k \eq \scalar{f}{f}^{1/2}$ for $f \in H^k(\G)$,
so $H^k(\G)$ forms a Hilbert space in its own right.

\item[8.]
\textit{Sobolev spaces of functions having homogeneous boundary values}.
The space $\TF$ of test functions can surely be endowed
with the norms $\abs{\cdot} _{k,p}$,
but this leads to uncomplete normed spaces.
Hence it is nearby to introduce the spaces $\WN{k,p}{\G}$ for $k \in \N$
by completion of the space $\TF$ {\wrt} the Sobolev norms:
\[
\WN{k,p}{\G} \, \eqdef \, \overline{\TF_{\abs{\cdot} _{k,p}}}.
\]
The elements of these spaces can be identified with functions in $\SobW$.
Furthermore,
the spaces $\WN{k,p}{\G}$ form closed linear subspaces
of the Sobolev spaces $\SobW$.
Clearly,
the spaces $\HN{k}{\G} \eqdef \HN{k,2}{\G}$ are Hilbert spaces
with inner products and norms as defined for $H^k(\G)$.
\end{itemize}

\subsubsection*{} %- examples of adjoint spaces
In the following we describe the adjoint spaces
for a couple of topological linear spaces $\banX$ arising in applied analysis
by quoting a linear homeomorphism $\banX^* \cong \banY$,
where $\banY$ has to be chosen in a suitable way.

\begin{itemize}[leftmargin=15pt,itemsep=6pt]
\item[1.] \textit{The adjoint space of $L^p(\G)$}.
Let the two real numbers $p,q >1$ be coupled
by the relation %$\frac{1}{p}+\frac{1}{q}=1$.
$p^{-1} \,{+}\, q^{-1} \eq 1$.
Then the duality statement $L^p(\G)^* \cong L^q(\G)$ for $1 < p < \infty$
is a consequence of the \textit{{\HOELDER} inequality},
which claims
that for $f \in L^p(\G)$ and $g \in L^q(\G)$
the product $f \cdot g$ is Lebesgue-integrable and satisfies the estimation
\[
\norm{f \cdot g}_1 \, \leq \, \norm{f}_p \cdot \norm{g}_q.
\]
Thus,
each function $f \in L^q(\G)$ gives rise to a continuous linear functional
$\phi_f$ on the space $L^p(\G)$ due to the mapping
\begin{equation}\label{LINEAR_FUNCTIONAL_ON_LP}
\phi_f(g) \eqdef \int_{\,\G} fg \, d\mu, \quad g \in L^p(\G).
\end{equation}
On the other hand,
there are no more continuous linear functionals on $L^p(\G)$
besides those defined by \eqref{LINEAR_FUNCTIONAL_ON_LP}.

\subparagraph*{} %- adjoint space of L1
In the special case $p \eq 1$ we obtain $L^1(\G)^* \cong L^{\infty}(\G)$.

\item[2.] \textit{The adjoint space of $C_c(\G)$.}
It is well-known from measure theory
that each continuous function $f\:\G \xto \R$ is measurable,
but $f$ is in general not Lebesgue-integrable.
If $f \in C_c(\G)$,
however,
then $f$ is bounded and this implies Lebesgue-integrability.

\subparagraph*{} %- Radon measures
A \textit{Radon measure},
defined on the Borel subsets of a domain $\G \subset \Rd$, 
is a positive Borel measure
which takes finite values on compact subsets of $\G$.
Letting $\mu$ be a Radon measure on $\G$,
the mapping $\phi_{\mu}\: C_c(\G) \xto \R$ defined by
\[
\phi_{\mu}(f) \eqdef \int_{\,\G} f \, d\mu
\]
forms a continuous linear functional,
since
\[
\int_{\,\G} f \, d\mu \, \leq \,
\norm{f}_{\infty} \, \cdot \, \mu(\supp{f}).
\]
In addition,
$\phi_{\mu}(f) \geq 0$ holds for all $f \in C_c(\G)$ with $f \geq 0$,
and a linear functional fulfilling this property is called \textit{positive}.
So every Borel measure $\mu$ gives rise to
a positive and continuous linear functional $\phi_{\mu}$ on $C_c(\G)$.

\subparagraph*{} %- Riesz-Markov theorem I
Conversely,
the question arises
whether each continuous linear functional $\phi \in C_c(\G)^*$
appears in this way.
The answer is given by the \textsc{Riesz-Markov} \textit{theorem},
quoting
that the representation of positive linear functionals on $\G$
can just be done by the Radon measures on $\G$:

\begin{theorem}\label{RIESZ-MARKOV_THEOREM}
Let $\phi$ be a positive and continuous linear functional
on the normed space $C_c(\G)$.
Then there is exactly one Radon measure
$\mu\:\Borel{\G} \xto [0,\infty]$ such that
\[
\phi(f) \eqdef \int_{\,\G} f \, d\mu, \quad \quad f \in C_c(\G).
\]
\end{theorem}

\item[3.] \textit{The adjoint space of $C_0(\G)$.}
A \textit{complex} measure $\mu \: \Borel{\G} \xto \C$
is said to be \textit{Radon},
if the so-called \textit{variation} $\abs{\mu} $
is Radon in the sense of a positive measure on $\G$.
Notice that $\abs{\mu} $ is always finite.
The space of complex Radon measures $\MCR{\G}$
proves to be complete {\wrt} the norm $\norm{\mu} := \abs{\mu} $
for $\mu \in \MCR{\G}$.

\paragraph*{} %- Riez-Markov theorem II
The following theorem of \textsc{Riesz-Markov} type
describes the adjoint space of $C_0(\G)$:

\begin{theorem}\label{RIESZ_MARKOV_THEOREM_II}
Let $\phi$ be a complex-valued and continuous linear functional
defined on the space $C_0(\G)$.
Then there is exactly one complex Radon measure $\mu$ such that
\[
\phi(f) \, \eqdef \, \int_{\,\G} f \, d\mu, \quad \quad f \in C_0(\G).
\]
\end{theorem}
In conclusion,
$C_0(\G)^* \cong \MCR{\G}$,
so $C_0(\G)^*$ is norm-isomorphic to the Banach space $\MCR{\G}$
of complex Radon measures.

\item[4.] \textit{The adjoint space of $C[a,b]$.}
We consider a further variant of the \textsc{Riesz-Markov} theorem
for a compact set $[a,b]$ in the real line.
Due to the properties of the Riemann-Stieltjes integral,
each function $g$ \textit{of bounded variation} on $[a,b]$,
gives rise to a continuous linear functional $F$ on the space $C[a,b]$
according to
\[
G(f) \, \eqdef \, \int_a^b f \, dg, \quad \quad (f \in C[a,b]).
\]
Let $BV(a,b)$ denote the linear space of all functions of bounded variation.
This space can be equipped with the norm
\[
\norm{f}_{bv} \, \eqdef \,
\abs{ f(a) } \, + \,
\sup_{P} \sum_{i=0}^{n_P-1} \abs{f(x_{i+1}-f(x_i)} ,
\]
where the supremum is taken over all finite partitions of $[a,b]$.
The space $BV(a,b)$ is actually a Banach space in this norm,
which is even norm-isomorphic to the space $\MCR{a,b}$
of Radon measures on $[a,b]$.

\subparagraph*{}
Now the \textsc{Riesz-Markov} theorem quotes that
\textit{every} continuous linear functional $G$ on $C[a,b]$ arises
as Riemann-Stieltjes integral
{\wrt} a certain function $g$ of bounded variation,
that is,
we have $C[a,b]^* \cong BV(a,b)$.

\item[5.] \textit{The adjoint space of $\D(\G)$.}
The notion of \textit{distribution} forms a generalization
of an ordinary function $f\:\G \xto \F$.
The corresponding subject in functional analysis,
called the \textit{theory of distributions},
was developed by \textsc{L.\,Schwartz} in the midth of the 20th century
to describe ideas and concepts of theoretical physics
by an exact mathematical formalism.
Distribution theory can be seen as a direct application
of the abstract theory of locally convex linear spaces
and their continuous linear functionals.

\subparagraph*{}
Remember the space $\D(\G)$ consisting of test functions and furnished
with a locally convex topology,
in which this space is complete but not metrizable.
Then every continuous linear functional $\disF\:\D(\G) \xto \C$ 
is called a \textit{Schwartz distribution}.
We denote $\D'(\G) \eqdef \D(\G)^*$ the adjoint space of $\D(\G)$
respectively the space of Schwartz distributions on $\G$.
It is also possible to characterize a Schwartz distribution
{\wrt} $\D$-convergence,
when applied to test functions:
a linear functional $\disF\:\TF \xto \C$
is a Schwartz distribution
iff $\disF(\phi_n) \xto 0$ for any $\D$-convergent sequence $(\phi_n)$
converges to the zero function.

\subparagraph*{} %- regular distributions
Notice that the linear functional $\disF_f$
derived from a locally integrable function $f$ by
\[
\disF_f(\phi) \, \eqdef \, \int_{\,\G} f \, \phi \, dx,
\quad \quad \phi \in \TF
\]
is continuous,
and this is also true for the linear functional
induced by the Radon measure $\mu$ according to
\[
\disF_{\mu}(\phi) \, \eqdef \, \int_{\,\G} \phi \, d\mu,
\quad \quad \, \phi \in \TF.
\]
Thus locally integrable functions as well as Radon measures
may be regarded as Schwartz distributions.
Notice that the space $\D'(\G)$ is indeed extremely extensive
and contains nearly all functions being important in applications,
amongst others the functions belonging to the well-known spaces
$\Lp$, $\Li$, $C(\G)$, $C(\clG)$ and $C_0(\G)$,
but also measures on $\G$,
which are regular enough in some sense.

\subparagraph*{}
A Schwartz distribution is said to be a \textit{regular distribution},
if it is the continuous linear functional $\disF_f$
induced by a locally integrable function $f \in \Lc$,
otherwise it is called \textit{singular}.
Well-known examples of singular distributions are
the \textit{Dirac distributions}
\[
\delta_a(\phi) \, \eqdef \, \phi(a)
\qfa \phi \in \TF
\quad \quad  (a \in \G).
\]

\item[6.] \textit{The adjoint space of $\E(\G)$.}
Regarding the {\FRECHET} space $\E(\G)$,
each element $\disF \in \E'(\G) \eqdef \E(\G)^*$ is called
a \textit{distribution with compact support}.
This notion is founded by the fact
that for such $\disF$ there is a compact set $K_{\disF} \subset \G$,
called the \textit{support} of $\disF$
such that
$\disF \phi \eq 0$ for all $\phi \in \E(\G)$
having support in $\G \setminus K_{\disF}$.
On the other hand,
each measurable function $f$ with compact support of course
gives rise to a compactly supported distribution $\disF_f \in \E'(\G)$.

\subparagraph*{}
Since it is always possible to increase an adjoint space
by decreasing the underlying base space,
there is a continuous inclusion $\E'(\G) \, \hookrightarrow \, \D'(\G)$,
which is dual to embedding
$\D(\G) \, \hookrightarrow \, \E(\G)$.
So we basically may regard
each compactly supported distribution as a Schwartz distribution.

\item[7.] \textit{The adjoint space of $\SRd$.}
Similar to the spaces $\D(\G)$,
we introduce the notion of \textit{tempered distribution}
as a continuous linear functional $\disT\:\SRd \xto \C$,
which generalizes that of bounded and slow-growing locally integrable functions.
The space of tempered distributions is usually denoted $\TD$.

\begin{comment}
\subparagraph*{}
The following criterion determines
under which conditions a linear functional $f$ on $\SRd$
forms a tempered distribution:
this is the case 
\iff
there is a constant $C > 0$ and integers $m,l \in \N$ such that
\[
\abs{ f(\phi) } \, \leq \,
C \cdot \sum_{\abs{\alpha} \leq l, \abs{\beta} \leq m} \, p_{\alpha,\beta}(\phi)
\qfa \phi \in \SRd\,,
\]
where the functions $p_{\alpha,\beta}$ are the seminorms
defined by \eqref{I_SRN_FAMILY_OF_SEMINORMS}.
\end{comment}

\subparagraph*{} %- comparison of distributional spaces
Comparing the known kinds of distributional spaces,
we obtain a chain of inclusions,
which is dual to that of the base spaces:
\[
\cal{E}'(\Rd) \, \subset \, \TD \, \subset \, \D'(\Rd).
\]
IF we endow each of these spaces with the \textit{weak*-topology},
then these embeddings prove to be continuous.
Hence a distribution having compact support can also be regarded
as a tempered distribution,
which in turn gives rise to a Schwartz distribution.

\item[8.] \textit{The adjoint space of $\HN{k}{\G}$}.
The Sobolev space \textit{with negative index} $H^{-k}(\G)$
is defined as the adjoint space
$H^{-k}(\G) \, \eqdef \, \HN{k}{\G}^*$.

\subparagraph*{}
Notice that not all elements of $H^{-k}(\G)$ can be represented
by regular functions $f\:\G \xto \F$,
but all belong to the space $\D'(\G)$ of Schwartz distributions.
In case of $k \eq 1$,
each $f \in H^{-1}(\G)$ is of the form
\[
f \= f_0 + \partial_1 f_1 + \ldots + \partial_d f_d,
\]
where $f_0,f_1,\ldots,f_d \in \Lh$.
There are similar representations for $k > 1$.

\subparagraph*{}
More generally,
the adjoint space of $\WN{k,p}{\G}$ ($1<p<\infty$) 
is the space $W^{-k,q}(\G)$,
where $p^{-1} \,{+}\, q^{-1} \eq 1$.
Like in the case $p \eq 2$,
the elements of $W^{-k,q}(G)$ can be regarded as Schwartz distributions on $G$.
\end{itemize}

\paragraph*{} %- orthonormal bases in Hilbert spaces L2(I)
Remember that
the Lebesgue spaces $L^2(X)$ are prototypes of separable Hilbert spaces,
where $X$ forms the base set of a certain measure space $(X,\Borel{X},\mu)$.

\subparagraph*{}
Our last subject is to enumerate orthonormal bases in $L^2(X)$,
where $X \eq (a,b) \subset \R$ is a finite or infinite open interval.

\begin{itemize}[leftmargin=15pt,itemsep=6pt]
\item[1.]
If $(a,b) \eq (-\pi,\pi)$,
the \textit{trigonometrical functions} given by the functions
\[
\frac{1}{\sqrt{2\pi}},         \;\,
\frac{1}{\sqrt{\pi}} \sin(x),  \;\,
\frac{1}{\sqrt{\pi}} \cos(x),  \;\,
\frac{1}{\sqrt{\pi}} \sin(2x), \;\,
\frac{1}{\sqrt{\pi}} \cos(2x), \;\, \dots
\]
form an orthonormal base in the Hilbert space $L^2((-\pi,\pi))$.

\item[2.]
In case of $(a,b)=(-1,1)$,
an orthonormal base, founded on polynomials,
is given by the so-called \textit{Legendre polynomials}
\[
P_n(x)
\, \eqdef \,
\frac{1}{2^n n!} \, \frac{d^n}{dx^n} (x^2-1)^n,
\quad n \in \N,
\]
together with the constant function $P_0(x) \eqdef 1$.
These functions form the only orthogonal system in $L^2(-1,1)$
that contains polynomials of any order.

\item[3.]
Let $(a,b)$ be the half-line $(0,\infty)$.
Then the sequence of functions
\[
L_n^{\alpha}(x)
\, \eqdef \,
x^{\alpha/2} e^{-x/2}
\cdot
\cal{L}_n^{\alpha}(x),
\quad n \in \N
\]
forms an orthonormal base in the space $L^2(0,\infty)$,
where $\cal{L}_n^{\alpha}$ are called the \textit{Laguerre polynomials},
defined by the equations
\[
\cal{L}_n^{\alpha}(x)
\, \eqdef \,
\frac{x^{-\alpha}e^x}{n!} \, \frac{d^n}{dx^n} (e^{-x}x^{n+\alpha}),
\quad n \in \N.
\]

\item[4.]
If $(a,b)$ is the entire set $\R$,
an orthonormal base of $L^2(\R)$ is given by the sequence
of \textit{Hermite functions}
\[
\cal{H}_n(x) \, \eqdef \, e^{-\frac{x^2}{2}} H_n(x),
\quad n \in \N,
\]
where $H_n$ are the \textit{Hermite polynomials}
\[
{H}_n(x)
\, \eqdef \,
(-1)^n \sum_{k_1+2k_2=n}
\frac{n!}{k_1! \, k_2!}
\,
(-1)^{k_1+k_2} \, (2x)^{k_1},
\quad n \in \N.
\]

\begin{comment}
\subparagraph*{} %- Hermite functions are eigenfunctions of the Plancherel oerator
The functions $\cal{H}_n$ belong
to the linear subspace $C^{\infty}(\R) \subset L^2(\R)$
and are the eigenfunctions of the Fourier-Plancherel operator $\FT_P$
{\wrt} the eigenvalues \-werte $+1,+i,-1,-i$,
distributed according to
\[
\begin{cases}[1.2]
\quad +1 &: \quad \cal{H}_0, \, \cal{H}_4, \, \cal{H}_8,    \; \, \ldots \\
\quad +i &: \quad \cal{H}_1, \, \cal{H}_5, \, \cal{H}_9,    \; \, \ldots \\
\quad -1 &: \quad \cal{H}_2, \, \cal{H}_6, \, \cal{H}_{10},    \, \ldots \\
\quad -i &: \quad \cal{H}_3, \, \cal{H}_7, \, \cal{H}_{11},    \, \ldots \\
\end{cases}
\]
\end{comment}
\end{itemize}

%------------------------------------------------------------------------------%
%                                 BIBLIOGRAPHY                                 %
%------------------------------------------------------------------------------%
%\newpage
%\thispagestyle{empty}
\vspace{10mm}

\titleformat{\section}
{\normalfont \bfseries \filcenter}
{\thesection.\;}{0mm}{}

\begin{thebibliography}{99}
\begin{small}
\bibitem{I_AF} R.\,Adams, J.\,Fournier:
\textit{Sobolev Spaces}, Academic Press Inc., 2003.
\bibitem{I_RS1} M.\,Reed, B.\,Simon:
\textit{Modern Methods of Mathematical Physics I: {\FA}}, Academic Press Inc., 1980.
\bibitem{I_Yo}  K.\,Yosida:
\textit{{\FA}}, 5.\,Edition, Springer Verlag, 1976.
\end{small}
\end{thebibliography}

%------------------------------------------------------------------------------%
%                               END OF DOCUMENT                                %
%------------------------------------------------------------------------------%
\end{document}
